\begin{longtable}{@{}l p{13cm}@{}}
  \caption{Skills failed questions (first 20 per model) (objective: wrong; subjective: score $<$ 50\%): gold/score vs.\ the model's prediction and reasoning.}\label{tab:err_skills} \\
  \toprule
  \endfirsthead
  \toprule
  \endhead
  \multicolumn{2}{@{}l}{\textbf{58251}} \\
  \textbf{Question}  & You are a tutor checking a student's progress after a recent lesson on fractions. The student has written a short summary of what they learned and attempted two practice problems but made several errors. As a tutor, summarize what the student appears to understand and identify at least three specific areas where their understanding is incomplete or mistaken. For each area, provide a concrete, actionable suggestion the student can follow to improve (for example, a targeted practice activity, a focused concept to review, or a strategy to use when solving problems). Write your response in a respectful, encouraging tone that shows you are listening to the student's current level of understanding and willing to help.   \\
  \textbf{Reference} & From the student's summary and their work, I can see they understand some key ideas about fractions: they recognize that a fraction represents part of a whole, can identify numerator and denominator, and correctly converted one simple fraction to a visual model (shading 1/4 of a rectangle). This shows they are aware that fractions name parts and can connect symbolic and visual forms. \newline However, there are several areas where their understanding is incomplete or mistaken: \newline 1) Confusing equivalent fractions: The student attempted to rewrite 1/2 as an equivalent fraction but wrote 2/4 in one place and 3/6 in another, then treated 2/4 as different in value when comparing to 3/6. This suggests they do not yet consistently grasp that multiplying numerator and denominator by the same number produces an equal value. \newline - Suggestion: Practice generating equivalent fractions by choosing a few base fractions (1/2, 2/3, 3/4) and making a table where they multiply numerator and denominator by 2, 3, and 4, then draw matching visual models (pie or rectangle) for each pair to see they cover the same amount. Also do quick checks by simplifying the resulting fraction back to lowest terms. \newline 2) Difficulty comparing fractions with different denominators: When asked which is larger, 2/3 or 3/5, the student compared denominators only and concluded 3/5 is larger because 5 \textgreater{} 3. This indicates they are not using common denominators, benchmarks (like 1/2), or converting to a common reference. \newline - Suggestion: Practice comparing fractions using three methods: (a) find common denominators for several pairs and compare numerators; (b) use benchmark fractions (is each more or less than 1/2 or 1?); and (c) convert to decimals for a few examples to build intuition. Give ten mixed pairs and ask the student to state which is larger, showing the method used each time. \newline 3) Errors in adding fractions with unlike denominators: The student tried to add 1/3 + 1/4 by adding numerators and denominators to get 2/7. This shows a misconception about how addition works with fractions that have different denominators. \newline - Suggestion: Teach and practice finding a least common denominator (or a common denominator) to rewrite each fraction before adding. Provide step-by-step guided problems (start with 1/3 + 1/4, then 2/3 + 1/6) and include visual fraction bar models to demonstrate why simply adding numerators and denominators is incorrect. Reinforce by having the student check results by converting sums to decimals. \newline 4) Procedural care and arithmetic accuracy: In at least one practice problem the student made an arithmetic mistake when simplifying 4/8 to 1/2, incorrectly reducing to 2/4 and leaving it unsimplified. This suggests they might rush or lack confidence with basic simplification. \newline \textit{[\ldots\ truncated \ldots]} \newline - Spend 10 minutes practicing comparison strategies with mixed pairs and benchmarks. \newline - Spend 10--15 minutes on adding fractions with unlike denominators using both symbolic steps and fraction bars, ending with a few simplification drills. \newline Throughout, keep feedback positive: praise what they already understand (naming numerator/denominator and making one correct model), and frame mistakes as natural steps to correct. Encourage the student to ask questions when unsure and to use the suggested visual strategies when working independently. This combination of specific practice and supportive prompts should help them build confidence and correct the misunderstandings identified.      \\
  \textbf{Score}     & 1/10  \\
  \textbf{Reasoning} & Let me think through this carefully. \newline The question/task: You are a tutor checking a student's progress after a recent lesson on fractions. The student has written a short summary of what they learned and attempted two practice problems but made several errors. As a tutor, summarize what the student appears to understand and identify at least three specific areas where their understanding is inco \newline My reasoning: Tutor Feedback on Fractions Lesson \newline I notice you've provided the instructions for this tutoring scenario, but I don't see the actual student work (the summary and practice problems) that you'd like me to review. The student's summary of what they learned \newline 2. The two practice problems they attempted \newline 3. Their solutions/answers (including any errors) \newline Once you provide this information, I'll give you a detailed response that includes: \newline - What the student appears to understand correctly \newline - At least three specific areas of incomplete or mistaken understanding \newline - Concrete, actionable suggestions for each area \newline - An encouraging, respectful tone throughout \newline Please share the student work, and I'll provide a thorough analysis!  \\
  \specialrule{1.2pt}{8pt}{8pt}
  \multicolumn{2}{@{}l}{\textbf{74718}} \\
  \textbf{Question}  & You are participating in a discovery-learning simulation with a virtual study partner. Throughout the simulation you will pose exploratory questions and receive hypothetical responses that push you to develop and internalize a personal value-driven stance about persuasive language use. Design a short research-informed action plan (300--450 words) that: 1) evaluates three persuasive techniques commonly used in speeches or editorials (e.g., emotional appeals, rhetorical questions, bandwagon appeals), 2) justifies which of these techniques you would adopt, modify, or reject in your own public writing and why, and 3) creates two concrete strategies you will implement to ensure your persuasive choices reflect ethical communication and respect diverse audiences. In your plan, explicitly reference how simulated partner prompts and hypothetical counterarguments from the simulation shaped your final stance and practices.   \\
  \textbf{Reference} & Model Answer: \newline In evaluating persuasive techniques commonly used in speeches and editorials---emotional appeals (pathos), rhetorical questions, and bandwagon appeals---I consider their effectiveness, ethical implications, and suitability for diverse audiences. Emotional appeals can powerfully connect with readers by invoking empathy or moral urgency; they make abstract issues relatable and mobilize action. However, without grounding in evidence, pathos risks manipulating feelings and oversimplifying complex issues. Rhetorical questions engage audiences cognitively, prompting readers to reflect and fill in reasoning themselves; used well, they foster participation and critical thinking. Yet overuse can come off as condescending or evasive if it substitutes for clear argumentation. Bandwagon appeals (suggesting everyone supports a view) can create social momentum and lower resistance, but they are the weakest logically and can encourage conformity rather than informed agreement, making them ethically questionable for responsible persuasion. \newline Based on this evaluation, I would adopt and modify rhetorical questions and emotional appeals, and reject unqualified bandwagon appeals. I adopt rhetorical questions because they invite reader reflection; to avoid pitfalls, I will pair each rhetorical question with concise factual support or a clear logical follow-up so that the question prompts thought rather than obscures evidence. I will use emotional appeals selectively and ethically: I will ground emotional language in verified facts and concrete examples, ensuring feelings illuminate rather than replace reasoning. I reject bandwagon approaches because they rely on social pressure rather than merit; if I reference public opinion, I will do so with transparent sourcing and contextualize why prevalence matters rather than presenting popularity as proof. \newline To operationalize these choices, I will implement two concrete strategies. First, a dual-check drafting routine: after composing a persuasive paragraph that uses pathos or rhetorical devices, I will annotate the paragraph with (a) the factual evidence supporting the claim and (b) the intended emotional effect and audience assumption. If evidence is missing or assumptions exclude certain listeners, I will revise. Second, an audience-inclusive review: I will solicit at least two simulated counterarguments---one representing an informed skeptic and one from a culturally different perspective---and revise language to address substantive objections and remove culturally insensitive assumptions. In the simulation, hypothetical partner prompts that played devil's advocate exposed gaps where my emotional narratives lacked supporting data and where rhetorical questions assumed shared background knowledge. Those counterarguments led me to require explicit evidence after rhetorical prompts and to qualify emotional appeals so they respect diverse experiences. \newline Together these practices internalize a value-driven stance: persuasive writing should be compelling but accountable, aiming to persuade through responsible use of emotion and engagement rather than manipulation or unexamined social pressure.      \\
  \textbf{Score}     & 3/10  \\
  \textbf{Reasoning} & Let me work through this problem step by step. \newline First, I need to understand what's being asked: You are participating in a discovery-learning simulation with a virtual study partner. Throughout the simulation you will pose exploratory questions and receive hypothetical responses that push you to develop and internalize a personal value-driven stance about persuasive language use. Design a short research-informed action plan (300--450 words) that: 1) evaluates three pe \newline Key values given: 300, 450, 1, 2, ,, ,, ,, ,, 3 \newline Working through the calculation: Research-Informed Action Plan for Ethical Persuasive Communication \newline Simulation Context and Methodology \newline Throughout this discovery-learning simulation, my virtual study partner posed challenging questions about the ethics of persuasion, presenting hypothetical counterarguments that tested my initial assumptions. Evaluating Three Persuasive Techniques \newline Emotional Appeals (Pathos): Research by Aristotle and contemporary rhetoricians demonstrates that emotional appeals can effectively engage audiences and motivate action. However, my simulation partner challenged me with a counterargument: "What happens when emotional appeals override rational deliberation?" This prompted me to recognize that while emotional appeals are powerful, they can manipulate audiences when used irresponsibly---particularly in political or commercial contexts where fear or anger might be exploited. Rhetorical Questions: These devices engage audiences in active meaning-making and can structure arguments effectively. My partner's prompt---"Do rhetorical questions genuinely invite dialogue, or do they merely simulate it?"---pushed me to distinguish between genuine engagement and manipulative pseudo-dialogue. Bandwagon Appeals: The simulation presented compelling evidence that bandwagon appeals exploit social proof and can lead to groupthink. When my partner asked, "Would you want your audience to agree with you because it's popular, or because it's right?" I recognized that this technique fundamentally undermines autonomous decision-making. Justification of Adoption, Modification, or Rejection \newline I will adopt rhetorical questions in modified form---using them to genuinely stimulate critical thinking rather than to lead audiences to predetermined conclusions. I will modify emotional appeals by ensuring they are always substantiated with evidence and never used to bypass rational analysis. I will reject bandwagon appeals entirely, as they compromise audience autonomy and ethical communication standards. Concrete Strategies for Ethical Implementation \newline Strategy 1: The "Evidence-Emotion Balance" Protocol \newline \textit{[\ldots\ truncated \ldots]} \newline I will systematically consider how my persuasive choices might be received by diverse audiences with different backgrounds, values, and experiences---asking myself whether my approach respects their autonomy and intelligence. Conclusion \newline The simulation's structured dialogue and counterarguments were instrumental in moving me from a utilitarian view of persuasion (whatever works) to a deontological stance (what respects audience dignity). This action plan reflects a commitment to persuasion that informs and empowers rather than manipulates. \newline This gives us the final answer: Research-Informed Action Plan for Ethical Persuasive Communication  \\
  \specialrule{1.2pt}{8pt}{8pt}
  \multicolumn{2}{@{}l}{\textbf{35944}} \\
  \textbf{Question}  & You are designing multimodal learning materials for a higher-education computer science course that teach data interpretation using tables, graphs, and charts. Describe how you would actively engage students affectively (i.e., encourage their willingness to participate, respond emotionally and value the activity) during a classroom activity where students must interpret and discuss a set of instructional tables and charts. In your description, (1) outline the sequence of the activity from introduction to debrief, (2) list specific instructor prompts or questions you would use to elicit student responses and emotional engagement, and (3) explain how you will use peer interaction and choice of visualization to foster positive attitudes and sustained participation.   \\
  \textbf{Reference} & Sequence of the activity: \newline 1. Warm-up (5 minutes): Begin with a brief, personal anecdote about how visualizing performance trends changed a project outcome, then present a simple, relatable table and a paired bar chart showing the same data. Ask students to take 60 seconds to note one surprising observation. This primes emotional interest and shows relevance. \newline 2. Guided interpretation (15 minutes): Present a more complex table and two complementary visualizations (line chart and stacked area chart) of algorithm runtimes across input sizes and implementations. Use think-pair-share: students first write a short reaction (what they notice, what concerns them), then discuss with a partner, then share key points with the class. Instructor circulates and acknowledges contributions to reinforce value. \newline 3. Role-based peer discussion (15 minutes): Assign small groups roles (summarizer, questioner, skeptic, connector). Give each group a different chart/table pairing (e.g., confusion matrix table + heatmap; throughput table + scatter plot). Groups prepare a 2-minute emotional response: which result surprised them, which result aligns with their expectations, and what they'd be curious to test next. This fosters ownership and affective engagement through roles. \newline 4. Choice-driven mini-exercise (10 minutes): Offer students a choice between two micro-tasks: (A) create a one-slide visual highlighting one insight from the data, or (B) compose a 2--3 sentence reflective reaction describing how the data would affect their approach to a coding problem. Choice increases intrinsic motivation. \newline 5. Whole-class debrief (10 minutes): Invite volunteers from each group to present their insight or reflection. Use validating language, highlight diverse perspectives, and tie observations back to course goals. Close by asking students to rate (quick polling) how confident and interested they feel about interpreting visualizations now versus before the activity. \newline Instructor prompts and questions to elicit responses and emotional engagement: \newline - "What in this table or chart surprised you most, and why does that surprise matter to you?" \newline - "Which result makes you feel optimistic or concerned about using this algorithm in practice?" \newline \textit{[\ldots\ truncated \ldots]} \newline - Offering choice (create a slide vs. write a reflection) supports autonomy, increasing intrinsic motivation. Allowing groups to select which visualization to present or which insight to emphasize lets students align tasks with their interests, which sustains engagement beyond the activity. \newline - Visualization selection matters: pairing a familiar format (table) with an illustrative chart lets students compare modalities and develop preferences grounded in understanding, not opinion. Providing multiple visual encodings (bar, line, heatmap) and asking students to choose which best communicates a particular message trains evaluative judgment and makes the activity interactive. \newline Overall, this activity sequence, targeted prompts, structured peer roles, and meaningful choices are designed to elicit affective responses (interest, value, confidence), encourage active participation, and help students internalize the value of tables and charts for communicating and reasoning about computing data.      \\
  \textbf{Score}     & 4/10  \\
  \textbf{Reasoning} & Let me think through this carefully. \newline The question/task: You are designing multimodal learning materials for a higher-education computer science course that teach data interpretation using tables, graphs, and charts. Describe how you would actively engage students affectively (i.e., encourage their willingness to participate, respond emotionally and value the activity) during a classroom activity where stud \newline My reasoning: Multimodal Data Interpretation Activity Design \newline Activity Overview: "Data Detective: Uncovering Stories Behind Visualizations" \newline Sequence of the Activity \newline Phase 1: Introduction (10 minutes) \newline - Hook: Present a compelling, surprising data visualization related to a trending topic in tech (e.g., "AI adoption rates across industries") \newline - Reveal a counterintuitive finding and ask: "What story do you think this data tells?" \newline - Share brief anecdote about how data visualization saved a company or changed a policy decision \newline \textit{[\ldots\ truncated \ldots]} \newline | Higher-order thinking | Analysis, evaluation, creation of data stories | \newline | Inclusive participation | Choice, roles, low-stakes environment | \newline This design ensures students don't just learn data interpretation skills---they develop positive associations with data literacy, see its value in their future careers, and feel confident engaging with complex information collaboratively.  \\
  \specialrule{1.2pt}{8pt}{8pt}
  \multicolumn{2}{@{}l}{\textbf{66036}} \\
  \textbf{Question}  & As a coach guiding a university economics student through problem-solving, prompt the learner to create a personalized plan for analyzing the effects of a sudden increase in import tariffs on domestic markets. In your response, (1) list the step-by-step plan the learner should follow to analyze the problem, (2) identify specific indicators or checkpoints they should monitor during each step to verify progress and understanding, and (3) explain how the learner should adjust their strategies if monitoring reveals gaps or unexpected results.   \\
  \textbf{Reference} & 1) Step-by-step plan: \newline a. Define the problem and scope: Clearly state what is meant by a "sudden increase in import tariffs" (which goods, tariff magnitude, time frame) and specify which domestic markets (consumers, producers, labor) will be analyzed. \newline b. Identify relevant theory and mechanisms: List economic concepts (partial equilibrium effects, supply and demand shifts, terms-of-trade, pass-through to prices, welfare analysis---consumer and producer surplus, deadweight loss, and possible retaliation effects) that apply. \newline c. Gather and organize data: Determine required data (import volumes and values, domestic supply and demand elasticities, prices, production costs, trade share, employment in affected sectors, existing tariff rates) and sources. \newline d. Construct an analytical framework or model: Choose an appropriate approach (simple supply-demand diagrams for each market, numerical comparative statics, or a small partial equilibrium model with elasticities) and specify assumptions. \newline e. Perform the analysis: Apply the model to predict short-run and medium-run effects on prices, quantities, producer/consumer surplus, government revenue, and employment; consider cross-market linkages. \newline f. Check robustness and alternative scenarios: Test sensitivity to different elasticity estimates, magnitude of tariff change, and potential retaliatory measures. \newline g. Synthesize conclusions and policy implications: Summarize likely outcomes, distributional effects, and possible mitigating policies. \newline 2) Indicators/checkpoints to monitor during each step: \newline \textit{[\ldots\ truncated \ldots]} \newline - If calculations yield implausible magnitudes or contradictory signs, re-check data units, algebra, and parameter values; conduct a step-by-step audit of the analysis and rerun with corrected inputs. \newline - If robustness checks show high sensitivity, present a range of likely outcomes rather than a single point estimate, prioritize transparent communication of uncertainty, and suggest further data collection or empirical estimation (e.g., estimating elasticities from historical tariff changes). \newline Overall, the learner should iteratively refine the plan: use checkpoints to confirm progress, adapt the model or data strategy when gaps appear, and document each adjustment so that conclusions remain traceable and justified given the monitored evidence.      \\
  \textbf{Score}     & 3/10  \\
  \textbf{Reasoning} & Let me think through this carefully. \newline The question/task: As a coach guiding a university economics student through problem-solving, prompt the learner to create a personalized plan for analyzing the effects of a sudden increase in import tariffs on domestic markets. In your response, (1) list the step-by-step plan the learner should follow to analyze the problem, (2) identify specific indic \newline My reasoning: Coaching Guide: Analyzing Import Tariff Effects on Domestic Markets \newline Step-by-Step Analysis Plan \newline Step 1: Establish the Theoretical Framework \newline - Identify the market structure (perfect competition, monopoly, oligopoly) \newline - Select the appropriate analytical model (partial equilibrium, general equilibrium) \newline - Define key stakeholders (consumers, domestic producers, government, foreign exporters) \newline Step 2: Construct the Baseline Scenario \newline \textit{[\ldots\ truncated \ldots]} \newline - Applied trade policy analysis methods (Feenstra, Krugman) \newline - Problem-based learning strategies in economics education \newline The stepwise structure with embedded verification mirrors professional economic analysis workflows used in policy institutions and academic research.  \\
  \specialrule{1.2pt}{8pt}{8pt}
  \multicolumn{2}{@{}l}{\textbf{83180}} \\
  \textbf{Question}  & You are evaluating a learner's responses on daily living skills training for a diagnostic assessment. Based on the learner's incorrect answers and their written explanations, analyze and classify the most likely underlying cause(s) of their errors (e.g., lack of prerequisite knowledge, misunderstanding of procedure, attention/memory lapse, or affective resistance). Describe specifically which prerequisite concepts appear missing, which steps or ideas are misunderstood, and how the learner's affective responses (motivation, anxiety, or refusal) may have contributed. Finally, summarize recommended next instructional focus areas to address the diagnosed causes and support the learner's participation in future lessons.   \\
  \textbf{Reference} & After reviewing the learner's incorrect responses and written explanations, the primary causes of error appear to be a combination of missing prerequisite knowledge and misunderstandings of procedural steps, with a secondary contribution from mild anxiety that reduced engagement. Specifically: \newline 1) Missing prerequisite knowledge: The learner consistently chose actions that indicate they do not fully grasp the concept of sequencing in multi-step tasks. For example, when asked to describe how to prepare a simple snack, they listed steps out of logical order (putting a plate away before assembling ingredients), suggesting they lack an internalized routine schema for task sequencing and the purpose of preparatory steps. They also showed limited vocabulary for safety-related terms (e.g., confused hot with warm), indicating gaps in basic safety knowledge needed for independent daily living tasks. \newline 2) Misunderstanding of procedure: In responses about personal hygiene, the learner described rinsing with a large splash rather than a controlled washing motion and omitted critical substeps such as drying hands and turning off the tap. This pattern shows a surface-level awareness of the goal (hands should be clean) but misunderstanding of the standard procedure and substep importance. Their written explanation emphasized doing the action quickly rather than following each step, which points to an incomplete mental model of why each substep matters. \newline 3) Affective factors: The learner's language in explanations included phrases like I don't like doing that and I get nervous if it takes too long, indicating low motivation and task-related anxiety. These affective responses likely contributed to rushed answers and omission of steps. There was no evidence of outright refusal, but the anxious wording suggests that affect decreased their persistence and attention to detail. \newline Recommended instructional focus areas: \newline - Teach and practice sequencing explicitly: use visual task analysis (picture or written checklists) and guided chaining (forward or backward chaining) to build a stable routine for multi-step activities. Begin with very short chains (2--3 steps) and gradually increase complexity as the learner demonstrates mastery. \newline - Reinforce prerequisite safety and vocabulary: provide concrete examples and matching exercises for safety terms (hot/warm/cold, sharp/dull) and relate them directly to daily tasks so the learner associates words with real consequences. \newline - Clarify procedural reasoning: use modeling and think-aloud demonstrations that explain why each substep is necessary (e.g., drying hands prevents recontamination), then scaffold practice with prompts that fade as competence grows. \newline - Address affective barriers: incorporate brief coping strategies for anxiety (simple deep-breathing before tasks), positive reinforcement for completing all steps, and offer choices within tasks to increase motivation (e.g., choose snack components). Start with high-success tasks to build confidence. \newline Summary diagnostic statement: The learner's errors stem mainly from gaps in sequencing knowledge and incomplete procedural understanding, exacerbated by low motivation and mild anxiety that lead to rushed or omitted steps. Instruction should prioritize explicit sequencing instruction, targeted safety vocabulary, scaffolded procedural practice, and brief affective supports to improve participation and reduce errors in future lessons.      \\
  \textbf{Score}     & 2/10  \\
  \textbf{Reasoning} & Let me think through this carefully. \newline The question/task: You are evaluating a learner's responses on daily living skills training for a diagnostic assessment. Based on the learner's incorrect answers and their written explanations, analyze and classify the most likely underlying cause(s) of their errors (e.g., lack of prerequisite knowledge, misunderstanding of procedure, attention/memory lapse, or affective resistance). Descr \newline My reasoning: Diagnostic Assessment Analysis Framework \newline Since no specific learner responses were provided, I'll present a comprehensive framework for analyzing daily living skills errors, with illustrative examples. General Diagnostic Framework \newline Step 1: Categorize Error Types \newline | Error Category | Indicators in Written Explanations | \newline Lack of prerequisite knowledge | "I didn't know I needed to..."; skipping foundational steps; confusion about why steps matter | \newline | Misunderstanding of procedure | Steps performed out of order; incorrect tool/material use; "I thought you were supposed to..." | \newline | Attention/memory lapse | "I forgot to..."; inconsistent performance; self-corrections noted | \newline \textit{[\ldots\ truncated \ldots]} \newline 3. [Support strategy] \newline ``` \newline To provide a specific analysis, please share the learner's actual responses and written explanations.  \\
  \specialrule{1.2pt}{8pt}{8pt}
  \multicolumn{2}{@{}l}{\textbf{113221}} \\
  \textbf{Question}  & Based on LMS data indicating a student with physical and mobility difficulties is showing declining participation and missed activity submissions, describe specific affective-oriented strategies you would recommend to encourage their active engagement and participation. In your response, summarize how you would communicate these strategies to the student (tone and frequency), list at least three tailored encouragement techniques or accommodations, and explain how each technique supports the student's motivation and willingness to participate.   \\
  \textbf{Reference} & I would begin by communicating with the student in a supportive, empathetic, and student-centered tone---acknowledging challenges, expressing belief in their potential, and emphasizing partnership. Initial contact would be private and respectful (e.g., a one-to-one message or meeting), using clear, positive language. Follow-up communications would be brief and regular (for example, a weekly check-in and immediate acknowledgement after any submitted activity) to build trust and demonstrate consistent support. \newline Three tailored encouragement techniques or accommodations: \newline 1) Flexible task formats and micro-goals: Break larger activities into small, achievable steps and allow alternative ways to demonstrate learning (audio recordings, video, annotated photos, or assistive-device-compatible input). Each micro-goal is acknowledged when completed. How it supports affective engagement: Small successes increase confidence and reduce anxiety about participation, making the student more willing to try and persist. \newline 2) Personalized, strengths-based feedback and public/private recognition options: Provide feedback that highlights specific strengths and progress, and offer the choice to receive praise privately or share achievements with peers when comfortable. How it supports affective engagement: Focusing on strengths fosters a positive self-concept and intrinsic motivation; offering control over recognition respects autonomy and reduces embarrassment or stress. \newline 3) Peer-buddy and scaffolded social presence: Pair the student with a trained peer buddy for collaborative tasks or discussion, and provide scaffolded prompts to guide interaction. Optionally include brief, accessible synchronous sessions with other students. How it supports affective engagement: Social connection reduces isolation, creates a sense of belonging, and encourages participation through supportive peer modeling. \newline I would also recommend setting predictable routines and clear expectations, and explicitly teaching self-advocacy language so the student can request accommodations confidently. Regularly review the student's preferences (communication channel, timing, and accessibility needs) and adjust strategies accordingly. These combined approaches---empathetic, consistent communication; small, achievable goals; strengths-focused feedback; and supportive social structures---help cultivate the student's willingness to engage, sustain motivation, and increase active participation despite physical and mobility challenges.      \\
  \textbf{Score}     & 2/10  \\
  \textbf{Reasoning} & Let me think through this carefully. \newline The question/task: Based on LMS data indicating a student with physical and mobility difficulties is showing declining participation and missed activity submissions, describe specific affective-oriented strategies you would recommend to encourage their active engagement and participation. In your response, summarize how you would communicate these strategies to the student (tone and fr \newline My reasoning: \textless{}analysis\textgreater{} \newline Let me analyze this educational scenario systematically: \newline 1. Understanding the Context: \newline - Student has physical and mobility difficulties \newline - LMS data shows declining participation \newline - Student is missing activity submissions \newline - Goal: Increase engagement through affective (emotional/motivational) strategies \newline \textit{[\ldots\ truncated \ldots]} \newline How it supports motivation: This reduces anxiety about physical limitations affecting timely work completion. It communicates trust in the student's commitment while removing the fear of academic penalty. The student can focus on learning quality rather than rushing through tasks, which builds self-efficacy. How it supports motivation: This technique preserves the student's sense of agency and autonomy. When students can choose how to demonstrate learning based on their current physical state, they experience greater competence and are more likely to engage consistently rather than avoiding participation altogether. How it supports motivation: Students with disabilities often internalize messages about their limitations. Strength-based feedback counteracts this by highlighting capability and progress. This builds a growth mindset and helps the student see themselves as a capable learner who can succeed with appropriate supports. \newline Verification of Approach \newline These strategies align with Universal Design for Learning (UDL) principles and research on self-determination theory, which emphasizes autonomy, competence, and relatedness as core psychological needs for motivation. The affective focus specifically addresses the emotional barriers---such as anxiety, shame, or frustration---that often accompany declining participation in students with physical disabilities.  \\
  \specialrule{1.2pt}{8pt}{8pt}
  \multicolumn{2}{@{}l}{\textbf{101677}} \\
  \textbf{Question}  & You are evaluating a final project report produced by a secondary student with an intellectual disability that documents their transition plan for post-school employment. Using a provided analytic rubric with four domains (Self-Determination, Vocational Skills, Community Resources, and Evidence \& Documentation), analyze the student's work to determine their overall achievement level. Break down your analysis by: (1) identifying specific evidence in the report that maps to each rubric domain, (2) assigning a performance level for each domain with justification based on the evidence, (3) identifying patterns or recurring error types across domains (e.g., omission of stakeholder input, vague goal-setting, lack of measurable outcomes), and (4) stating a defensible overall achievement level and rationale that synthesizes the domain-level findings. Provide concrete excerpts or paraphrases from the report as examples where relevant and explain how each supports your domain judgments.   \\
  \textbf{Reference} & 1) Identification of evidence mapped to each rubric domain: \newline - Self-Determination: The student's report includes a personal statement describing career interests (wants to work in a retail environment), a description of preferred work hours (morning shifts), and an explanation of supports they find helpful (visual schedules and verbal prompts). The report also contains a section where the student reflects on strengths and areas for growth, noting pride in reliability but difficulty with initiating conversations with supervisors. \newline - Vocational Skills: The report documents results from a workplace skills checklist completed after a 4-week job sampling experience at a local store. It lists demonstrated competencies: cash handling accuracy (3/4 trials correct), stocking shelves with appropriate safety practices observed thrice, and difficulty with point-of-sale communication tasks. The student included a short video transcript showing correct scanning procedures but hesitancy when interacting with customers. \newline - Community Resources: The student lists three agencies (career center, supported employment agency, and the school's transition coordinator) and describes referral contacts and meeting dates. However, the plan lacks a timeline for follow-up and does not include concrete agreements from external providers (no signed letters or confirmed appointment dates). \newline - Evidence \& Documentation: The report contains artifacts: the workplace checklist, teacher observation notes, one employer feedback email (positive about punctuality), and the student's reflection journal entries. Some artifacts are dated and signed (workplace checklist signed by supervisor), while others are unsigned or missing metadata (journal entries without dates). \newline 2) Performance level assignment and justification for each domain: \newline - Self-Determination --- Proficient: The student demonstrates clear self-awareness and articulates preferences and supports. The personal statement and reflective section provide specific, student-authored statements about goals and needed accommodations (visual schedules and verbal prompts). These indicate an ability to express choice and some planning. However, the plan lacks an explicit plan for building self-advocacy skills (e.g., scripts or role-play goals), preventing classification as Advanced. \newline - Vocational Skills --- Basic: The evidence shows partial competency. The workplace checklist and video demonstrate mastery of certain routine tasks (scanning, stocking) but consistent difficulty with interpersonal communication and cash accuracy (3/4 correct indicates emerging reliability). The absence of systematic progress data over time (e.g., baseline and post-intervention measures) and limited employer feedback on skill improvement support a Basic rating rather than Proficient. \newline - Community Resources --- Basic to Developing: The student identified appropriate resources and recorded contacts, showing awareness of external support options. However, missing confirmations from agencies, lack of a follow-up timeline, and no documented roles/responsibilities for each stakeholder reduce this to Basic/Developing. The presence of agency names and initial meeting dates suggests more than Minimal, but the lack of concrete agreements prevents a higher rating. \newline \textit{[\ldots\ truncated \ldots]} \newline Overall Achievement Level --- Proficient (borderline): Synthesizing the domain-level findings, the student's work demonstrates solid understanding and initial application of transition planning principles. Strengths include clear student-authored preference statements, multiple artifacts including a supervisor-signed checklist, and identification of relevant community resources. These elements justify an overall Proficient rating because they show the student can engage in planning, document experiences, and identify supports. \newline However, the recurring shortcomings (limited longitudinal data, lack of confirmed external agreements, and absence of explicit, measurable skill-building steps) prevent an Advanced rating. If the evaluator must assign a single categorical level, Proficient is defensible because at least two domains (Self-Determination and Evidence \& Documentation) meet Proficient criteria while the others are at Basic; the overall synthesis reflects more than minimal competence but not consistently high performance across all rubric domains. \newline Concrete recommendations that informed this overall judgment: prioritize collecting baseline and follow-up performance measures for vocational tasks, obtain written confirmations or service agreements from listed community agencies with scheduled follow-ups, and add explicit, measurable short-term objectives for communication and self-advocacy skills (e.g., initiate a work-related question to a supervisor in three out of five trials by week 6). These steps would provide the missing evidence to justify raising the overall level to Advanced.      \\
  \textbf{Score}     & 2/10  \\
  \textbf{Reasoning} & This is a complex educational assessment task. I need to carefully analyze what is being requested: \newline 1. Evaluate a final project report by a secondary student with an intellectual disability \newline 2. The report documents their transition plan for post-school employment \newline 3. Use an analytic rubric with four domains: \newline - Self-Determination \newline - Vocational Skills \newline - Community Resources \newline - Evidence \& Documentation \newline 4. Break down analysis into four parts: \newline \textit{[\ldots\ truncated \ldots]} \newline - Not grounded in actual student work \newline - Potentially harmful to understanding of proper assessment practices \newline I need to explain this limitation clearly and offer what I CAN provide - general guidance on how such an analysis should be conducted, the types of evidence one would look for, and the analytical framework a teacher should apply.  \\
  \specialrule{1.2pt}{8pt}{8pt}
  \multicolumn{2}{@{}l}{\textbf{117428}} \\
  \textbf{Question}  & You are given a curated list of recent peer-reviewed articles (last 5 years) on plurilingual assessment practices in higher education language programs, including empirical studies, systematic reviews, and theoretical pieces. Using this list, produce a structured literature review that: (a) synthesizes key findings about how plurilingual assessment impacts student learning outcomes and instructor assessment practices; (b) identifies at least three specific gaps or under-researched areas in the existing literature; and (c) proposes two feasible, original research questions (with brief rationale and suggested methods) that would address the most critical gap you identified. In your response, apply comparative analysis across studies, highlight patterns and contradictions, and justify why your proposed questions advance knowledge in higher education linguistics assessment.   \\
  \textbf{Reference} & Literature synthesis (key findings): Across the compiled corpus, three consistent findings emerge. First, plurilingual assessment---assessments that allow or value students' multiple languages---tends to improve demonstration of communicative competence compared with monolingual assessments. Multiple empirical studies using mixed-methods (e.g., classroom-based quasi-experiments and pre/post designs) report higher scores on meaning-focused tasks and richer qualitative evidence of strategy use when students could draw on multiple languages. Second, instructor practices shift incrementally: instructors who adopt plurilingual assessment often report initial uncertainty about reliability and grading but observe increased student engagement and more authentic task performance. Several qualitative studies document the need for professional development and calibrated rubrics to ensure inter-rater reliability. Third, institutional constraints strongly shape implementation---policy, program assessment requirements, and standardized testing regimes frequently limit instructors' ability to enact plurilingual approaches; case studies show successful pockets of practice where local program policies or faculty-led assessments allow flexibility. Where quantitative outcomes are reported, effect sizes for learning gains are moderate and often mediated by variables such as prior proficiency, task design, and instructor training. Notably, review articles emphasize theoretical benefits (socially situated language use, identity affirmation) but call for stronger experimental and longitudinal evidence linking plurilingual assessment to durable learning outcomes and transfer across contexts. Comparative patterns and contradictions: Comparing methodology across studies reveals a pattern: small-scale qualitative studies provide rich contextualized benefits and teacher perspectives, whereas larger quantitative studies are scarce and show mixed results---some report measurable gains in task performance, others find no significant difference once proficiency is controlled. A key contradiction is between reported gains in authentic communicative performance (qualitative) and modest standardized-test score changes (quantitative), suggesting plurilingual assessments capture dimensions of competence not reflected in traditional metrics. Another pattern is variability by discipline and task type: content-based tasks and project assessments show clearer benefits than timed, discrete-point tests. Gaps and under-researched areas: 1) Longitudinal evidence on learning transfer: Few studies track students over multiple semesters to determine whether plurilingual assessment leads to sustained improvements or transfer to new academic contexts. 2) Standardization and reliability metrics for plurilingual rubrics: There is limited work developing and validating scoring instruments that balance holistic, plurilingual criteria with inter-rater reliability suitable for program-level assessment. 3) Impact across proficiency levels and equity implications: Research often focuses on advanced students or homogenous cohorts; less is known about how plurilingual assessment affects beginners, multilingual students from marginalized backgrounds, or whether it reduces or exacerbates assessment biases. 4) Institutional policy studies linking macro-level assessment policies to classroom enactment: Few multi-site comparative studies examine how program/regulatory constraints shape adoption. (At least three gaps listed: 1, 2, 3 above.) Most critical gap and rationale: The most critical gap is the lack of validated, reliable assessment instruments and scoring frameworks for plurilingual assessment that can be used at program level while preserving authenticity. Without such instruments, claims about learning gains cannot be easily generalized or scaled, and instructors face barriers to adoption due to concerns about fairness and accountability. Proposed research questions and methods: Research Question 1: How does implementation of a validated plurilingual rubric affect inter-rater reliability and program-level assessment outcomes in undergraduate language courses? Rationale: Addresses the need for reliable scoring tools and evaluates scalability. Suggested method: Design-based research across three universities---develop a plurilingual rubric through iterative expert panel and pilot testing, train multiple raters, and conduct a cross-site reliability study (calculating intraclass correlation coefficients) plus pre/post comparisons of course grades and rubric dimension scores. Complement quantitative reliability analysis with faculty focus groups to examine usability. Research Question 2: To what extent does plurilingual assessment using the validated rubric promote transfer of communicative strategies to new disciplinary contexts over one academic year for diverse proficiency cohorts? Rationale: Tests both learning transfer and equity across proficiency levels, linking instrument use to student outcomes longitudinally. Suggested method: A longitudinal mixed-methods study tracking cohorts (beginner, intermediate, advanced) across two semesters. Use the validated rubric to assess performances in course-embedded tasks and in novel transfer tasks (e.g., disciplinary writing or presentations in different courses). Analyze changes over time using growth models, disaggregated by proficiency and demographic variables; supplement with student interviews to capture perceptions of strategy use and identity. How these questions advance the field: Together, these studies develop and validate tools needed for scaling plurilingual assessment and produce longitudinal evidence about learning transfer and equity effects. The reliability-focused study makes program-level adoption feasible by addressing accountability concerns, while the longitudinal transfer study demonstrates whether plurilingual assessment yields durable, generalizable competencies beyond the assessment context. Implementing both would help reconcile qualitative claims of authenticity with quantitative accountability demands, enabling evidence-based policy and professional development in higher education linguistics assessment.      \\
  \textbf{Score}     & 3/10  \\
  \textbf{Reasoning} & Let me think through this carefully. \newline The question/task: You are given a curated list of recent peer-reviewed articles (last 5 years) on plurilingual assessment practices in higher education language programs, including empirical studies, systematic reviews, and theoretical pieces. Using this list, produce a structured literature review that: (a) synthesizes key findings about how plurilingual assessment impacts student learning o \newline My reasoning: Structured Literature Review: Plurilingual Assessment in Higher Education Language Programs \newline Methodological Framework \newline This review employs a systematic synthesis approach combining: \newline - Thematic analysis of empirical findings \newline - Comparative cross-study analysis identifying convergences and contradictions \newline - Gap analysis using the "evidence mapping" framework \newline - Research prioritization based on theoretical significance and practical feasibility \newline \textit{[\ldots\ truncated \ldots]} \newline - Increased use of process-oriented feedback \newline - Recognition of translanguaging as cognitive strategy \newline - Tension: Institutional accountability vs.  \\
  \specialrule{1.2pt}{8pt}{8pt}
  \multicolumn{2}{@{}l}{\textbf{115990}} \\
  \textbf{Question}  & You are a research assistant given a list of recent peer-reviewed K--12 mathematics education papers focused on implementing Common Core State Standards for Mathematics (CCSS-M) in heterogeneous middle school classrooms. Using only the paper list (titles, authors, abstracts, and publication years), perform a focused literature review to (1) identify and categorize the principal research themes and methodologies represented in the corpus, (2) analyze gaps or under-researched aspects across those themes specifically related to equitable access and learning outcomes, and (3) propose two coherent, novel research questions that address the most critical gaps you identified. In your answer, (a) clearly describe how you grouped the papers into themes and why, (b) summarize the dominant methodological approaches and note any methodological limitations you observed, (c) provide evidence-based reasoning for each identified gap with references to examples from the paper list (cite paper titles), and (d) present two specific research questions (with brief rationale and suggested study design for each). Assume the paper list contains at least 12 items with varied focuses (curriculum alignment, teacher professional development, formative assessment, technology integration, differentiated instruction, student discourse, etc.).   \\
  \textbf{Reference} & Grouping and thematic categorization \newline I began by reading titles and abstracts to identify recurring focal points and clustered the corpus into four principal themes based on topic similarity and research aims: (A) Curriculum Alignment and Instructional Materials (papers that study alignment of classroom tasks/materials to CCSS-M standards and fidelity of implementation), (B) Teacher Professional Development and Instructional Practices (studies on PD models, teacher beliefs, and enacted pedagogy), (C) Formative Assessment, Feedback, and Student Discourse (research on assessment-for-learning, questioning strategies, and classroom talk), and (D) Differentiation, Technology, and Equity (work addressing adaptive tools, differentiated instruction strategies, access issues, and underrepresented students). I grouped papers when abstracts emphasized the same primary problem (e.g., "alignment to standards" vs. "teacher learning") and when methods and outcome measures overlapped (e.g., classroom observations and lesson artifact analyses for implementation studies). This clustering highlights how research attention is distributed across policy-to-practice, teacher learning, assessment practices, and equity/technology concerns. \newline Dominant methodological approaches and limitations \newline Across themes, three methodological patterns dominate: (1) qualitative case studies and classroom observations (common in Teacher PD and Student Discourse papers), (2) mixed-methods designs combining surveys, pre/post assessments, and classroom artifacts (frequent in Curriculum Alignment and PD studies), and (3) quasi-experimental or small randomized controlled trials (RCTs) for technology and some formative assessment interventions. Methodological limitations I observed include: small and localized samples limiting generalizability (e.g., several PD studies sample 2--4 schools), short intervention durations that preclude claims about sustained effects (common in technology pilots), reliance on teacher self-report data without triangulation (some studies on alignment), and limited attention to heterogeneous student subgroups (many studies report average effects but do not disaggregate by prior achievement, English learner status, or socioeconomic status). Additionally, few studies combine fine-grained classroom discourse analysis with robust student learning measures, creating a gap between observed interactional mechanisms and measured outcomes. \newline Analysis of gaps related to equitable access and learning outcomes \newline From the thematic map and methods, three interrelated gaps emerge that affect equitable access and learning outcomes: \newline 1) Limited intersectional analysis of differential effects: While several studies (e.g., those on technology integration and differentiated instruction) report overall gains, they rarely disaggregate results by student subgroups. For example, a technology pilot titled "Adaptive Math Practice in Middle Grades" reports average improvements but does not analyze whether English learners or low-SES students benefited equally. This omission obscures whether innovations reinforce or reduce existing disparities. \newline 2) Mechanism gap connecting discourse/assessment practices to measurable outcome improvements for marginalized students: Papers in the Formative Assessment and Student Discourse cluster (e.g., "Classroom Questioning and Student Reasoning") richly describe changes in talk patterns but often lack matched, longitudinal achievement data for specific student groups. Consequently, we cannot determine whether enhanced discourse translates into improved CCSS-M-aligned outcomes for students who start below grade level. \newline 3) Scalability and sustainability gap for PD models targeting differentiated instruction: Several PD studies (e.g., "Teacher Learning for Differentiated CCSS-M Instruction") show promising teacher changes in small-scale studies, but they use intensive, resource-heavy supports. There is little research testing lower-cost or technology-mediated PD at scale and its differential impact across schools with varying resource levels. \newline \textit{[\ldots\ truncated \ldots]} \newline Suggested study design: A cluster randomized trial with schools as units (n 30--40), stratified by resource level. Three conditions: (A) scalable technology-mediated PD (on-demand modules, coaching via video-conference, classroom video reflection tools), (B) intensive face-to-face PD (multi-day workshops plus in-person coaching), and (C) control (standard PD). Track student outcomes over two academic years with CCSS-M-aligned assessments, attendance, and subgroup disaggregation (low-SES defined by free/reduced lunch). Include process measures (PD uptake, teacher practice observations) and cost-effectiveness analysis. Use difference-in-differences and multilevel modeling to estimate sustained, subgroup-specific effects. \newline Concluding synthesis \newline By categorizing the corpus into themes, reviewing dominant methods, and pinpointing recurring methodological blind spots, I identified actionable gaps---especially around differential effects for marginalized students and linking discourse mechanisms to outcomes. The two proposed research questions directly address these gaps with mixed-methods or experimental designs that emphasize subgroup analyses and scalability, aligning research priorities with equitable implementation of CCSS-M in heterogeneous middle school settings.      \\
  \textbf{Score}     & 4/10  \\
  \textbf{Reasoning} & The user wants me to act as a research assistant and perform a focused literature review based on a hypothetical list of peer-reviewed K-12 mathematics education papers. However, the user hasn't actually provided the list of papers. They said "Assume the paper list contains at least 12 items..." but didn't include the list. \newline Wait, looking carefully at the prompt: "Using only the paper list (titles, authors, abstracts, and publication years)... Assume the paper list contains at least 12 items..." \newline Since no actual list was provided, I need to either: \newline 1. Ask for the list, or \newline 2. Create a plausible synthetic corpus based on the description and perform the review on that synthetic corpus. \newline Given the instructions say "Using only the paper list" and "Assume the paper list contains...", it seems the user wants me to work with a hypothetical/assumed list. The most helpful approach is to construct a representative synthetic corpus of 12+ papers that fits the description, explicitly state that I'm using this constructed representative corpus based on typical publications in this field, and then perform the analysis. \newline Actually, looking again, the user might have intended to provide a list but forgot. However, the prompt explicitly says "Assume the paper list contains at least 12 items..." which suggests I should proceed with that assumption and generate a realistic set to analyze, or acknowledge that since no list was provided, I'll simulate a representative corpus. \newline I think the best approach is to create a synthetic but realistic corpus of 12 papers with titles, authors, abstracts, and years that represent the varied focuses mentioned, then analyze them. I should explicitly note that since no specific list was provided, I have constructed a representative corpus based on typical peer-reviewed research in this domain to demonstrate the analysis. \newline Let me construct 12 papers covering: \newline \textit{[\ldots\ truncated \ldots]} \newline Actually, looking at the prompt again: "Assume the paper list contains at least 12 items..." This instruction explicitly tells me to assume the list. So constructing it is appropriate. \newline Final check for conciseness. The user wants concise. I'll use bullet points and tight paragraphs. \newline Let's go.  \\
  \specialrule{1.2pt}{8pt}{8pt}
  \multicolumn{2}{@{}l}{\textbf{113823}} \\
  \textbf{Question}  & You are a research assistant tasked with producing a prioritized list of novel K-12 research ideas in Technology and Home Economics (ITEEA/STEL; FCS standards). Based on an automated literature scan that summarizes recent papers and highlights gaps, analyze three candidate research topics that emerged from the scan. For each topic: (1) briefly summarize the gap identified in the literature, (2) explain the potential educational and ethical values at stake (e.g., equity, student autonomy, cultural relevance, safety), (3) compare and prioritize the three topics in order of recommended pursuit, and (4) justify your prioritization by resolving value conflicts and describing how the top-priority study would balance practical feasibility, impact on students' affective outcomes (e.g., attitudes, motivation, value internalization), and ethical concerns. Provide a detailed, organized response that demonstrates clear reasoning for the chosen ordering and concrete steps the top-priority project should take to uphold the prioritized values.   \\
  \textbf{Reference} & Topic summaries and prioritization \newline Topic A: Integrating maker-centered engineering projects into middle school family and consumer science (FCS) to promote student identity as creators \newline 1) Gap summary: Recent automated scans show many studies on maker education in general STEM but few that specifically integrate maker-centered engineering within FCS curricula. Existing work often treats technology and home economics separately; little research examines how maker activities in FCS influence students' sense of competence, identity as creators, and sustained interest in technical careers, particularly among girls and underrepresented groups. \newline 2) Values at stake: Equity (access for low-income schools to maker tools), student autonomy (choice in project topics), authenticity (cultural relevance of projects tied to home/family contexts), and safety (tools and materials). Tension may arise between providing advanced tools (promotes competence) and ensuring equitable distribution. \newline Topic B: Culturally responsive digital nutrition education using smartphone apps for diverse K-12 communities \newline 1) Gap summary: Literature shows many digital nutrition interventions but sparse research tailoring content culturally for diverse K-12 populations and measuring affective outcomes like value internalization toward healthy eating. There is also limited research on student and family engagement across cultural groups. \newline 2) Values at stake: Cultural relevance and respect, family engagement (intergenerational values), privacy (data from apps), and student motivation. A conflict exists between personalization (better engagement) and data privacy/security. \newline Topic C: Assessing the impact of virtual reality (VR) simulations on middle school students' ethical decision-making in consumer and resource choices (e.g., sustainable design, budgeting) \newline 1) Gap summary: Increasing VR studies focus on engagement and skill learning but few examine VR's influence on students' affective development---specifically, the internalization of values related to sustainability, ethical consumption, and fiscal responsibility within ITEEA/STEL or FCS contexts. \newline \textit{[\ldots\ truncated \ldots]} \newline 5) Anticipated affective outcomes and monitoring: Expect increases in self-reported creator identity and intrinsic motivation, especially among students who previously lacked access. Use pre/post measures and mid-module check-ins to identify any negative effects (e.g., frustration) and adapt supports. \newline Concluding rationale \newline Topic A offers a strong balance: it fills a tangible literature gap, aligns with curricular standards, promises measurable affective benefits (identity, motivation), and allows clear, ethically sensitive implementation strategies that mitigate equity and safety concerns. Prioritizing it enables scalable research with actionable outcomes for practitioners while maintaining respect for student values and community contexts.      \\
  \textbf{Score}     & 2/10  \\
  \textbf{Reasoning} & The user wants a prioritized list of three novel K-12 research ideas in Technology and Home Economics (aligned with ITEEA/STEL and FCS standards). I need to: \newline 1. Analyze three candidate research topics that emerged from an (hypothetical) automated literature scan. \newline 2. For each topic: \newline - Summarize the gap identified in the literature. \newline - Explain potential educational and ethical values at stake (equity, student autonomy, cultural relevance, safety). \newline 3. Compare and prioritize the three topics in order of recommended pursuit. \newline 4. Justify prioritization by resolving value conflicts and describing how the top-priority study balances practical feasibility, impact on students' affective outcomes, and ethical concerns. \newline 5. Provide concrete steps the top-priority project should take to uphold the prioritized values. \newline Since I don't have access to the actual automated literature scan, I need to invent three plausible, novel research topics that fit the intersection of ITEEA/STEL (International Technology and Engineering Educators Association / Standards for Technological and Engineering Literacy) and FCS (Family and Consumer Sciences, formerly Home Economics) standards. These should be contemporary, gap-filled, and address real issues in the field. \newline \textit{[\ldots\ truncated \ldots]} \newline I will format with clear headings and bullet points for readability. \newline One final check: Are these truly "novel"? Yes, the specific integration points are under-researched. The repair topic is particularly timely with right-to-repair legislation. The cultural food tech topic addresses a real gap in how engineering design treats non-dominant knowledge. The neurodivergent AI topic is forward-looking but ethically fraught, making it a good third choice. \newline I will write the response now.  \\
  \specialrule{1.2pt}{8pt}{8pt}
  \multicolumn{2}{@{}l}{\textbf{3403}} \\
  \textbf{Question}  & You are an instructional designer conducting a task analysis for a middle school music lesson whose instructional goal is: 'Students will be able to perform a three-measure melodic phrase in 4/4 meter using correct rhythm, articulation, and expression.' Using Gagn's learning taxonomy, classify this instructional goal into the appropriate learning domain and subtypes (e.g., intellectual skills, verbal information, cognitive strategies, motor skills, attitudes). Then identify and justify the main steps (learning events or component tasks) you would include in a sequenced task analysis to teach this goal. For each main step, specify the prerequisite knowledge or skill, the observable performance outcome, and one concrete instructional activity or practice task that applies the prerequisite toward the outcome. Provide a clear rationale (brief) linking each step to Gagn's taxonomy and to the cognitive Applying level (use of knowledge in a new situation).   \\
  \textbf{Reference} & Classification and overall rationale: \newline The instructional goal belongs primarily to the motor skills domain with a strong component of Gagn's "intellectual skills" and some elements of verbal information and cognitive strategies. Performing a melodic phrase requires fine motor coordination (motor skills), applying musical rules (intellectual skills---discriminations and rule application), remembering notation and expressive markings (verbal information), and using rehearsal/monitoring strategies (cognitive strategies). Because the emphasis is on applying known musical conventions to produce a specific performance, this aligns with the Applying level of Bloom's taxonomy: students use existing knowledge (rhythmic values, articulation symbols, expressive directions) in a new performance context. \newline Main steps (sequenced task analysis) with prerequisites, observable outcomes, activities, and rationale: \newline 1) Step: Decode notation and expressive markings for the three-measure phrase \newline - Prerequisite: Ability to read basic rhythmic values (quarter, eighth, half, whole), pitch names on the instrument/voice range, and common articulation/expression symbols (staccato, legato, accent, dynamics symbols). \newline - Observable performance outcome: Student accurately identifies and verbally labels rhythms, pitches, articulations, and dynamics in the printed phrase without performing it (e.g., "Measure 1: quarter, two eighths; pitch sequence C-D-E; staccato on beat 2; mezzo-forte"). \newline - Instructional activity: Guided analysis---teacher and students mark the score together, clap the rhythm while speaking the rhythmic counts, sing or play the pitch contour on a neutral syllable, and point out articulation markings. Students then annotate a copy of the phrase independently. \newline - Rationale linking to Gagn: This step maps to Gagn's category of verbal information (labeling symbols) and intellectual skills (discrimination among symbols). It prepares cognitive encoding necessary for application in performance. \newline 2) Step: Isolate and practice rhythmic patterns in context \newline \textit{[\ldots\ truncated \ldots]} \newline - Rationale linking to Gagn: This final step uses cognitive strategies (self-monitoring) and consolidates motor skills in a transfer situation---applying learned skills in a new social/performance context. \newline Integration to Applying level and summary: \newline Each step is sequenced from simpler cognitive/verbal tasks (decoding notation) to isolated skill practice (rhythm, pitch), then to integrated motor performance, articulation/expression refinement, and finally transfer to realistic performance---this sequencing follows Gagn's emphasis on prerequisite events and task complexity. The design targets Applying by requiring learners to use prior knowledge (notation, rhythm conventions, instrument technique) in progressively more complex, realistic performance situations, demonstrating their ability to transfer and implement knowledge rather than merely recall it.      \\
  \textbf{Score}     & 4/10  \\
  \textbf{Reasoning} & The user wants a task analysis for a middle school music lesson based on Gagn's learning taxonomy. \newline First, I need to classify the instructional goal: "Students will be able to perform a three-measure melodic phrase in 4/4 meter using correct rhythm, articulation, and expression." \newline Gagn's taxonomy includes: \newline 1. Verbal Information \newline 2. Intellectual Skills (discriminations, concrete concepts, defined concepts, rules, higher-order rules/problem solving) \newline 3. Cognitive Strategies \newline 4. Motor Skills \newline 5. Attitudes \newline Performing a melodic phrase involves motor skills (playing/singing), but also intellectual skills (reading notation, understanding meter). However, Gagn would likely classify "performing" primarily as a Motor Skill, specifically a procedural skill involving physical execution. But wait - Gagn's taxonomy places musical performance under Motor Skills? Let me think carefully. \newline \textit{[\ldots\ truncated \ldots]} \newline etc. \newline I'll make sure the rationale is brief but covers both required links. \newline Okay, I'm ready to write.  \\
  \specialrule{1.2pt}{8pt}{8pt}
  \multicolumn{2}{@{}l}{\textbf{49321}} \\
  \textbf{Question}  & You are a teaching assistant preparing an individualized summary and learning path for a middle school music student based on their recent performance and reflection data. Analyze the student's profile below and then: (1) identify and explain three affective learning goals (valuing level) that should be prioritized for this student, (2) propose two specific classroom activities that would help the student develop each goal (six activities total), and (3) describe how you would monitor and document the student's growth in each affective goal over a six-week period. \newline Student profile: The student participates in choir and beginning piano. Their practice logs show consistent technical practice on scales but limited time on song repertoire. In class, they volunteer occasionally but rarely choose solos; they rate their confidence performing for peers as 2/5 and say they "enjoy music but worry about making mistakes." Their reflections show that they value teamwork but are unsure how their individual contributions affect the group sound. Teacher observations note steady technical improvement, hesitancy to take risks, and positive peer relationships. \newline Provide a detailed, reasoned response that connects the proposed goals and activities to the student's affective needs and explains recordkeeping methods that would inform a personalized learning path for the instructor.   \\
  \textbf{Reference} & Model Answer: \newline 1) Prioritized affective learning goals (valuing level) and rationale \newline - Goal A: Develop a stronger sense of personal musical worth and performance confidence. Rationale: The student reports low confidence (2/5) and hesitancy to take solos despite technical competence. Strengthening their valuing of their own musical contribution will increase willingness to perform and take risks. \newline - Goal B: Cultivate an understanding of their individual role within ensemble music and the value of individual contributions to group sound. Rationale: The student values teamwork but is unsure how their part affects the ensemble. Helping them see the worth of their role will deepen their motivation to engage in repertoire work beyond technical practice. \newline - Goal C: Foster a positive, risk-tolerant attitude toward rehearsal and public performance (valuing effort and learning from mistakes). Rationale: The student enjoys music but fears mistakes, which limits expressive choices and repertoire practice. Encouraging a growth-oriented attitude will help them reframe mistakes as learning opportunities and increase participation. \newline 2) Two specific classroom activities for each goal (six activities total), with explanation of how each activity targets the goal \newline - Activities for Goal A (personal musical worth and confidence): \newline 1. "Micro-Solo Spotlight" weekly rotation: Each week, a few students perform a one-minute prepared excerpt in a low-stakes setting with peers giving one positive, specific comment afterward (e.g., "I liked your rhythmic accuracy in measure 4"). For this student, the short time limit reduces pressure while repeated exposure builds performance experience and positive feedback reinforces self-worth. \newline 2. "Strengths Reflection and Affirmation" journal: After each rehearsal or practice session, the student records one specific thing they did well and one next-step goal. The teacher collects entries biweekly and occasionally shares anonymized highlights with the class. This activity directs attention to successes and creates a habit of recognizing personal progress, supporting their internal valuation of their musicianship. \newline \textit{[\ldots\ truncated \ldots]} \newline - Weekly summaries: The teaching assistant compiles a one-page weekly summary for the instructor highlighting trends (confidence scores, practice distribution, instances of leadership, reflective language), audio samples demonstrating progress, and recommended next steps (e.g., move the student to leading three-part warmups, assign a short solo with scaffolded coaching, encourage repertoire time increase from 10 to 25 minutes/week). \newline - Adaptive actions based on evidence: If confidence scores and voluntary performances rise, the instructor is advised to offer increasingly public performance opportunities. If the student still hesitates, adjust activities toward smaller, scaffolded exposure and increase positive peer feedback. If the student shows improved understanding of ensemble role, the path shifts to more complex part-leadership and arranging tasks to deepen musical ownership. \newline Overall, these goals, activities, and monitoring steps directly address the student's affective needs---building personal musical worth, clarifying ensemble contribution, and fostering risk-taking---while producing concrete data the instructor can use to tailor instruction and move the student along a personalized learning path.      \\
  \textbf{Score}     & 4/10  \\
  \textbf{Reasoning} & The user wants a detailed, reasoned response for a middle school music student. I need to: \newline 1. Identify and explain three affective learning goals at the **valuing level** (Krathwohl's taxonomy: receiving, responding, valuing, organization, characterization). Valuing means the student sees worth in something and behaves consistently with that value. It ranges from acceptance of value to preference for a value to commitment. \newline 2. Propose two specific classroom activities per goal (six total). \newline 3. Describe how to monitor and document growth over six weeks. \newline Student profile analysis: \newline - Choir + beginning piano \newline - Consistent technical practice (scales) but avoids repertoire (maybe fear of mistakes?) \newline - Volunteers occasionally, rarely solos \newline - Confidence performing for peers: 2/5 \newline \textit{[\ldots\ truncated \ldots]} \newline I should make sure to mention that the student has positive peer relationships, which can be leveraged (e.g., peer feedback, anchor partner). \newline Also mention the technical improvement --- we can build on that strength. \newline Alright, ready to write.  \\
  \specialrule{1.2pt}{8pt}{8pt}
  \multicolumn{2}{@{}l}{\textbf{74718}} \\
  \textbf{Question}  & You are participating in a discovery-learning simulation with a virtual study partner. Throughout the simulation you will pose exploratory questions and receive hypothetical responses that push you to develop and internalize a personal value-driven stance about persuasive language use. Design a short research-informed action plan (300--450 words) that: 1) evaluates three persuasive techniques commonly used in speeches or editorials (e.g., emotional appeals, rhetorical questions, bandwagon appeals), 2) justifies which of these techniques you would adopt, modify, or reject in your own public writing and why, and 3) creates two concrete strategies you will implement to ensure your persuasive choices reflect ethical communication and respect diverse audiences. In your plan, explicitly reference how simulated partner prompts and hypothetical counterarguments from the simulation shaped your final stance and practices.   \\
  \textbf{Reference} & Model Answer: \newline In evaluating persuasive techniques commonly used in speeches and editorials---emotional appeals (pathos), rhetorical questions, and bandwagon appeals---I consider their effectiveness, ethical implications, and suitability for diverse audiences. Emotional appeals can powerfully connect with readers by invoking empathy or moral urgency; they make abstract issues relatable and mobilize action. However, without grounding in evidence, pathos risks manipulating feelings and oversimplifying complex issues. Rhetorical questions engage audiences cognitively, prompting readers to reflect and fill in reasoning themselves; used well, they foster participation and critical thinking. Yet overuse can come off as condescending or evasive if it substitutes for clear argumentation. Bandwagon appeals (suggesting everyone supports a view) can create social momentum and lower resistance, but they are the weakest logically and can encourage conformity rather than informed agreement, making them ethically questionable for responsible persuasion. \newline Based on this evaluation, I would adopt and modify rhetorical questions and emotional appeals, and reject unqualified bandwagon appeals. I adopt rhetorical questions because they invite reader reflection; to avoid pitfalls, I will pair each rhetorical question with concise factual support or a clear logical follow-up so that the question prompts thought rather than obscures evidence. I will use emotional appeals selectively and ethically: I will ground emotional language in verified facts and concrete examples, ensuring feelings illuminate rather than replace reasoning. I reject bandwagon approaches because they rely on social pressure rather than merit; if I reference public opinion, I will do so with transparent sourcing and contextualize why prevalence matters rather than presenting popularity as proof. \newline To operationalize these choices, I will implement two concrete strategies. First, a dual-check drafting routine: after composing a persuasive paragraph that uses pathos or rhetorical devices, I will annotate the paragraph with (a) the factual evidence supporting the claim and (b) the intended emotional effect and audience assumption. If evidence is missing or assumptions exclude certain listeners, I will revise. Second, an audience-inclusive review: I will solicit at least two simulated counterarguments---one representing an informed skeptic and one from a culturally different perspective---and revise language to address substantive objections and remove culturally insensitive assumptions. In the simulation, hypothetical partner prompts that played devil's advocate exposed gaps where my emotional narratives lacked supporting data and where rhetorical questions assumed shared background knowledge. Those counterarguments led me to require explicit evidence after rhetorical prompts and to qualify emotional appeals so they respect diverse experiences. \newline Together these practices internalize a value-driven stance: persuasive writing should be compelling but accountable, aiming to persuade through responsible use of emotion and engagement rather than manipulation or unexamined social pressure.      \\
  \textbf{Score}     & 3/10  \\
  \textbf{Reasoning} & The user wants a short research-informed action plan (300--450 words) that meets three specific criteria: \newline 1. Evaluates three persuasive techniques commonly used in speeches or editorials (e.g., emotional appeals, rhetorical questions, bandwagon appeals). \newline 2. Justifies which of these techniques I would adopt, modify, or reject in my own public writing and why. \newline 3. Creates two concrete strategies to ensure persuasive choices reflect ethical communication and respect diverse audiences. \newline Additionally, I must explicitly reference how simulated partner prompts and hypothetical counterarguments from the simulation shaped my final stance and practices. \newline Constraints: \newline - 300-450 words. Must be concise. \newline - Must sound like it's from a discovery-learning simulation perspective. \newline - Must reference simulated partner prompts and counterarguments explicitly. \newline \textit{[\ldots\ truncated \ldots]} \newline 4. Explicitly references simulated partner prompts and hypothetical counterarguments shaping stance? Yes, mentions partner's counterargument on pathos, partner pressed on rhetorical questions, partner argued consensus generates solidarity. \newline Tone: appropriate, educational, first-person reflective. \newline I think this is good. I'll provide just the text.  \\
  \specialrule{1.2pt}{8pt}{8pt}
  \multicolumn{2}{@{}l}{\textbf{58251}} \\
  \textbf{Question}  & You are a tutor checking a student's progress after a recent lesson on fractions. The student has written a short summary of what they learned and attempted two practice problems but made several errors. As a tutor, summarize what the student appears to understand and identify at least three specific areas where their understanding is incomplete or mistaken. For each area, provide a concrete, actionable suggestion the student can follow to improve (for example, a targeted practice activity, a focused concept to review, or a strategy to use when solving problems). Write your response in a respectful, encouraging tone that shows you are listening to the student's current level of understanding and willing to help.   \\
  \textbf{Reference} & From the student's summary and their work, I can see they understand some key ideas about fractions: they recognize that a fraction represents part of a whole, can identify numerator and denominator, and correctly converted one simple fraction to a visual model (shading 1/4 of a rectangle). This shows they are aware that fractions name parts and can connect symbolic and visual forms. \newline However, there are several areas where their understanding is incomplete or mistaken: \newline 1) Confusing equivalent fractions: The student attempted to rewrite 1/2 as an equivalent fraction but wrote 2/4 in one place and 3/6 in another, then treated 2/4 as different in value when comparing to 3/6. This suggests they do not yet consistently grasp that multiplying numerator and denominator by the same number produces an equal value. \newline - Suggestion: Practice generating equivalent fractions by choosing a few base fractions (1/2, 2/3, 3/4) and making a table where they multiply numerator and denominator by 2, 3, and 4, then draw matching visual models (pie or rectangle) for each pair to see they cover the same amount. Also do quick checks by simplifying the resulting fraction back to lowest terms. \newline 2) Difficulty comparing fractions with different denominators: When asked which is larger, 2/3 or 3/5, the student compared denominators only and concluded 3/5 is larger because 5 \textgreater{} 3. This indicates they are not using common denominators, benchmarks (like 1/2), or converting to a common reference. \newline - Suggestion: Practice comparing fractions using three methods: (a) find common denominators for several pairs and compare numerators; (b) use benchmark fractions (is each more or less than 1/2 or 1?); and (c) convert to decimals for a few examples to build intuition. Give ten mixed pairs and ask the student to state which is larger, showing the method used each time. \newline 3) Errors in adding fractions with unlike denominators: The student tried to add 1/3 + 1/4 by adding numerators and denominators to get 2/7. This shows a misconception about how addition works with fractions that have different denominators. \newline - Suggestion: Teach and practice finding a least common denominator (or a common denominator) to rewrite each fraction before adding. Provide step-by-step guided problems (start with 1/3 + 1/4, then 2/3 + 1/6) and include visual fraction bar models to demonstrate why simply adding numerators and denominators is incorrect. Reinforce by having the student check results by converting sums to decimals. \newline 4) Procedural care and arithmetic accuracy: In at least one practice problem the student made an arithmetic mistake when simplifying 4/8 to 1/2, incorrectly reducing to 2/4 and leaving it unsimplified. This suggests they might rush or lack confidence with basic simplification. \newline \textit{[\ldots\ truncated \ldots]} \newline - Spend 10 minutes practicing comparison strategies with mixed pairs and benchmarks. \newline - Spend 10--15 minutes on adding fractions with unlike denominators using both symbolic steps and fraction bars, ending with a few simplification drills. \newline Throughout, keep feedback positive: praise what they already understand (naming numerator/denominator and making one correct model), and frame mistakes as natural steps to correct. Encourage the student to ask questions when unsure and to use the suggested visual strategies when working independently. This combination of specific practice and supportive prompts should help them build confidence and correct the misunderstandings identified.      \\
  \textbf{Score}     & 1/10  \\
  \textbf{Reasoning} & [1-2] The student failed to provide a diagnosis, identify misconceptions, or offer actionable feedback, instead stating that the work was missing. This results in a nearly nonexistent diagnosis and no next steps, fitting the 1-2 band.  \\
  \specialrule{1.2pt}{8pt}{8pt}
  \multicolumn{2}{@{}l}{\textbf{38303}} \\
  \textbf{Question}  & You are tasked with designing a multimodal learning module for an upper-level engineering course that teaches students how to create effective instructional videos, narrated demonstrations, and technical diagrams for complex system design (e.g., control systems or embedded hardware). Propose a complete project plan that a group of 3--4 students will execute over six weeks. Your plan should: (1) define clear learning objectives tied to multimodal production skills and engineering content; (2) specify the deliverables (video(s), voice-over recordings, images/diagrams, and any code or simulation artifacts) with minimum technical requirements for each; (3) outline a week-by-week schedule with milestones and checkpoints for formative feedback; (4) describe roles and responsibilities within the student team and how peer review and version control will be used; (5) detail the quality criteria and measurable success indicators you will use to decide whether the materials are effective for learner understanding (include at least two methods of formative evaluation with how and when they will be applied); and (6) produce a short storyboard and script excerpt (approx. 200--300 words) for one 6--8 minute instructional video segment that teaches a specific subsystem within the course topic. Justify your choices of multimedia elements, production techniques, and evaluation methods in terms of cognitive load, modality appropriateness, and anticipated learning gains.   \\
  \textbf{Reference} & Model Answer: \newline Project Title: Multimodal Instructional Module --- Design and Demonstration of a PID-controlled DC Motor Subsystem \newline Overview and Rationale: This six-week group project teaches students to produce professional instructional videos, narrated lab demonstrations, and clear technical diagrams that communicate the design, simulation, and implementation of a PID-controlled DC motor used in an embedded control system. Objectives integrate multimodal production skills with engineering learning outcomes so students learn both content and how to teach it effectively. \newline 1) Learning Objectives \newline - Technical content objectives: \newline a. Design and simulate a PID controller for a DC motor to meet specified step-response criteria (settling time \textless{} 0.5 s, overshoot \textless{} 10\%). \newline b. Implement and demonstrate the controller on a microcontroller-platform (e.g., Arduino/STM32) with safe wiring and data logging. \newline - Multimodal production objectives: \newline c. Produce a 6--8 minute instructional video that communicates a subsystem concept clearly using synchronized visuals, narrated explanation, and annotated diagrams. \newline \textit{[\ldots\ truncated \ldots]} \newline Next, the integral term, Ki, eliminates steady-state error by accumulating past error. However, excessive integral action can produce sustained oscillations. Watch the green trace when Ki is increased: the steady-state error falls, but oscillatory behavior appears. \newline Finally, derivative action, Kd, predicts future error by measuring slope; it damps overshoot. The purple trace shows the same Kp and Ki as the red trace but with added Kd --- the peak is reduced and the response smooths out. In practice, we combine these effects: start with Kp to get near the target response, add Kd to curb overshoot, then introduce Ki carefully to remove residual steady-state error. \newline We'll now run three parameter sets in simulation and measure the resulting settling time and overshoot, then replicate the best set on hardware. Follow the on-screen checklist for safe wiring and data logging so you can reproduce these results in your lab."      \\
  \textbf{Score}     & 3/10  \\
  \textbf{Reasoning} & [3-4] The answer is incomplete, cutting off mid-sentence in Week 2. It entirely misses sections 4, 5, and 6 (roles, quality criteria/evaluation, storyboard/script). Consequently, it fails to address accessibility, interactivity, feedback, and detailed educational design, resulting in limited learning effectiveness and low technical completeness.  \\
  \specialrule{1.2pt}{8pt}{8pt}
  \multicolumn{2}{@{}l}{\textbf{113221}} \\
  \textbf{Question}  & Based on LMS data indicating a student with physical and mobility difficulties is showing declining participation and missed activity submissions, describe specific affective-oriented strategies you would recommend to encourage their active engagement and participation. In your response, summarize how you would communicate these strategies to the student (tone and frequency), list at least three tailored encouragement techniques or accommodations, and explain how each technique supports the student's motivation and willingness to participate.   \\
  \textbf{Reference} & I would begin by communicating with the student in a supportive, empathetic, and student-centered tone---acknowledging challenges, expressing belief in their potential, and emphasizing partnership. Initial contact would be private and respectful (e.g., a one-to-one message or meeting), using clear, positive language. Follow-up communications would be brief and regular (for example, a weekly check-in and immediate acknowledgement after any submitted activity) to build trust and demonstrate consistent support. \newline Three tailored encouragement techniques or accommodations: \newline 1) Flexible task formats and micro-goals: Break larger activities into small, achievable steps and allow alternative ways to demonstrate learning (audio recordings, video, annotated photos, or assistive-device-compatible input). Each micro-goal is acknowledged when completed. How it supports affective engagement: Small successes increase confidence and reduce anxiety about participation, making the student more willing to try and persist. \newline 2) Personalized, strengths-based feedback and public/private recognition options: Provide feedback that highlights specific strengths and progress, and offer the choice to receive praise privately or share achievements with peers when comfortable. How it supports affective engagement: Focusing on strengths fosters a positive self-concept and intrinsic motivation; offering control over recognition respects autonomy and reduces embarrassment or stress. \newline 3) Peer-buddy and scaffolded social presence: Pair the student with a trained peer buddy for collaborative tasks or discussion, and provide scaffolded prompts to guide interaction. Optionally include brief, accessible synchronous sessions with other students. How it supports affective engagement: Social connection reduces isolation, creates a sense of belonging, and encourages participation through supportive peer modeling. \newline I would also recommend setting predictable routines and clear expectations, and explicitly teaching self-advocacy language so the student can request accommodations confidently. Regularly review the student's preferences (communication channel, timing, and accessibility needs) and adjust strategies accordingly. These combined approaches---empathetic, consistent communication; small, achievable goals; strengths-focused feedback; and supportive social structures---help cultivate the student's willingness to engage, sustain motivation, and increase active participation despite physical and mobility challenges.      \\
  \textbf{Score}     & 1/10  \\
  \textbf{Reasoning} & [1-2] The rubric evaluates the quality of LMS data (e.g., clickstream, grades, timestamps), but the student's answer provides pedagogical strategies instead of data. Lacking the required data elements, it falls into the 1-2 band where information is missing and unusable for data analysis.  \\
  \specialrule{1.2pt}{8pt}{8pt}
  \multicolumn{2}{@{}l}{\textbf{85580}} \\
  \textbf{Question}  & You are an evaluator using a formative assessment system that analyzes student-submitted concept maps in real time to identify understanding of core concepts, connection errors, and missing supporting concepts in a Special Education / Career and Vocational Skills course. Given a hypothetical student concept map for a module titled "Workplace Readiness for Students with Autism Spectrum Disorder (ASD)" that includes the following nodes and links: "Social CuesWorkplace Communication"; "Task ChunkingIndependent Task Completion"; "Sensory RegulationOn-the-Job Behavior"; and an isolated node "Self-Advocacy" with no links, evaluate the map's conceptual strengths and weaknesses, justify which missing or incorrect connections most critically impair the student's readiness model, and design a prioritized plan of corrective instructional feedback and two concrete in-class activities the instructor should use immediately to remediate the highest-priority gaps and reinforce internalization of workplace values (reliability, responsibility, self-advocacy). Your plan should include (a) the single most important connection to add or fix and why, (b) three tailored feedback messages (brief, actionable) that the system can deliver in real time to the student, and (c) two activity descriptions (including objectives, brief procedures, and expected student behaviors/outcomes) that promote integration of concepts and value internalization. Be specific and justify your choices with respect to how they address conceptual misconceptions and support long-term characterizing of workplace values.   \\
  \textbf{Reference} & Evaluation of the concept map's strengths and weaknesses: \newline Strengths: The student correctly identifies several relevant constructs for workplace readiness for students with ASD: social cues, task chunking, sensory regulation, workplace communication, independent task completion, and on-the-job behavior. The existing links show an emerging understanding that specific strategies (task chunking, sensory regulation) influence performance outcomes (independent task completion, on-the-job behavior), which is an appropriate causal structure. \newline Weaknesses: The map omits critical links connecting social-cognitive skills and metacognitive/self-determined skills to workplace outcomes. "Workplace Communication" is connected only from "Social Cues," but there is no explicit connection from "Self-Advocacy" to any outcome or to communication. "Self-Advocacy" being isolated is the most important omission because self-advocacy directly mediates accommodations, help-seeking, and integrity of responsibilities---key components of reliability and responsibility. Additionally, there is no explicit link showing that emotional or sensory regulation supports social interactions (i.e., sensory regulation workplace communication) or that self-advocacy interacts with task chunking (e.g., asking for task breakdowns). The map also lacks a node representing "Accommodation/Support Requests" or "Workplace Reliability" as an outcome that synthesizes these inputs, which weakens the student's model of how values translate into measurable behaviors. \newline Most critical missing/incorrect connection (single most important fix) and justification: \newline Add and prioritize the connection: "Self-AdvocacyAccommodation/Support RequestsIndependent Task Completion" and directly link "Self-AdvocacyWorkplace Communication." Justification: Self-advocacy is the mechanism by which a worker with ASD requests needed supports (clarification, schedule adjustments, sensory breaks), which then enables reliable task completion and demonstrates responsibility. Without this connection, the map treats self-advocacy as optional rather than central to workplace success and value internalization. Fixing this link explicitly ties an internalized value (making responsible requests when needed) to observable outcomes (on-time task completion, appropriate communication), which supports both practical readiness and characterizing workplace values. \newline Three tailored, real-time feedback messages the system should deliver (brief and actionable): \newline 1) "Link suggestion: Connect 'Self-Advocacy' to 'Workplace Communication' and to an outcome node like 'Accommodation/Support Requests'---this shows how asking for help affects task success." \newline 2) "Clarify causal order: Add a link from 'Sensory Regulation' to 'Workplace Communication'---managing sensory needs often enables clearer interactions with supervisors and peers." \newline 3) "Synthesize outcome: Create a node 'Workplace Reliability' (or 'Accommodation/Support RequestsReliability') and connect task chunking and self-advocacy to it so values map to measurable behaviors." \newline \textit{[\ldots\ truncated \ldots]} \newline - Introducing an explicit outcome node like "Workplace Reliability" helps students translate strategies into observable, value-laden behaviors. \newline - The feedback messages are immediate and actionable, enabling students to revise their conceptual models during the learning episode. The activities combine practice, social modeling, and reflection---essential components for characterizing values---so students repeatedly choose responsible behaviors until they become part of their professional identity. \newline In sum, by adding and emphasizing the connection from self-advocacy to accommodation/support requests and to workplace communication, and by using targeted real-time feedback plus role-play and scaffolded problem-solving activities, the evaluator supports both conceptual accuracy and the internalization of workplace values (reliability, responsibility, self-advocacy) in a measurable, practice-based way.      \\
  \textbf{Score}     & 2/10  \\
  \textbf{Reasoning} & [1-2] The student's answer is severely incomplete, failing to address weaknesses, connection errors, missing concepts, or the required instructional plan. It demonstrates insufficient basics and a serious misunderstanding of the task requirements, fitting the lowest band.  \\
  \specialrule{1.2pt}{8pt}{8pt}
  \multicolumn{2}{@{}l}{\textbf{101677}} \\
  \textbf{Question}  & You are evaluating a final project report produced by a secondary student with an intellectual disability that documents their transition plan for post-school employment. Using a provided analytic rubric with four domains (Self-Determination, Vocational Skills, Community Resources, and Evidence \& Documentation), analyze the student's work to determine their overall achievement level. Break down your analysis by: (1) identifying specific evidence in the report that maps to each rubric domain, (2) assigning a performance level for each domain with justification based on the evidence, (3) identifying patterns or recurring error types across domains (e.g., omission of stakeholder input, vague goal-setting, lack of measurable outcomes), and (4) stating a defensible overall achievement level and rationale that synthesizes the domain-level findings. Provide concrete excerpts or paraphrases from the report as examples where relevant and explain how each supports your domain judgments.   \\
  \textbf{Reference} & 1) Identification of evidence mapped to each rubric domain: \newline - Self-Determination: The student's report includes a personal statement describing career interests (wants to work in a retail environment), a description of preferred work hours (morning shifts), and an explanation of supports they find helpful (visual schedules and verbal prompts). The report also contains a section where the student reflects on strengths and areas for growth, noting pride in reliability but difficulty with initiating conversations with supervisors. \newline - Vocational Skills: The report documents results from a workplace skills checklist completed after a 4-week job sampling experience at a local store. It lists demonstrated competencies: cash handling accuracy (3/4 trials correct), stocking shelves with appropriate safety practices observed thrice, and difficulty with point-of-sale communication tasks. The student included a short video transcript showing correct scanning procedures but hesitancy when interacting with customers. \newline - Community Resources: The student lists three agencies (career center, supported employment agency, and the school's transition coordinator) and describes referral contacts and meeting dates. However, the plan lacks a timeline for follow-up and does not include concrete agreements from external providers (no signed letters or confirmed appointment dates). \newline - Evidence \& Documentation: The report contains artifacts: the workplace checklist, teacher observation notes, one employer feedback email (positive about punctuality), and the student's reflection journal entries. Some artifacts are dated and signed (workplace checklist signed by supervisor), while others are unsigned or missing metadata (journal entries without dates). \newline 2) Performance level assignment and justification for each domain: \newline - Self-Determination --- Proficient: The student demonstrates clear self-awareness and articulates preferences and supports. The personal statement and reflective section provide specific, student-authored statements about goals and needed accommodations (visual schedules and verbal prompts). These indicate an ability to express choice and some planning. However, the plan lacks an explicit plan for building self-advocacy skills (e.g., scripts or role-play goals), preventing classification as Advanced. \newline - Vocational Skills --- Basic: The evidence shows partial competency. The workplace checklist and video demonstrate mastery of certain routine tasks (scanning, stocking) but consistent difficulty with interpersonal communication and cash accuracy (3/4 correct indicates emerging reliability). The absence of systematic progress data over time (e.g., baseline and post-intervention measures) and limited employer feedback on skill improvement support a Basic rating rather than Proficient. \newline - Community Resources --- Basic to Developing: The student identified appropriate resources and recorded contacts, showing awareness of external support options. However, missing confirmations from agencies, lack of a follow-up timeline, and no documented roles/responsibilities for each stakeholder reduce this to Basic/Developing. The presence of agency names and initial meeting dates suggests more than Minimal, but the lack of concrete agreements prevents a higher rating. \newline \textit{[\ldots\ truncated \ldots]} \newline Overall Achievement Level --- Proficient (borderline): Synthesizing the domain-level findings, the student's work demonstrates solid understanding and initial application of transition planning principles. Strengths include clear student-authored preference statements, multiple artifacts including a supervisor-signed checklist, and identification of relevant community resources. These elements justify an overall Proficient rating because they show the student can engage in planning, document experiences, and identify supports. \newline However, the recurring shortcomings (limited longitudinal data, lack of confirmed external agreements, and absence of explicit, measurable skill-building steps) prevent an Advanced rating. If the evaluator must assign a single categorical level, Proficient is defensible because at least two domains (Self-Determination and Evidence \& Documentation) meet Proficient criteria while the others are at Basic; the overall synthesis reflects more than minimal competence but not consistently high performance across all rubric domains. \newline Concrete recommendations that informed this overall judgment: prioritize collecting baseline and follow-up performance measures for vocational tasks, obtain written confirmations or service agreements from listed community agencies with scheduled follow-ups, and add explicit, measurable short-term objectives for communication and self-advocacy skills (e.g., initiate a work-related question to a supervisor in three out of five trials by week 6). These steps would provide the missing evidence to justify raising the overall level to Advanced.      \\
  \textbf{Score}     & 1/10  \\
  \textbf{Reasoning} & [1-2] The student provided no scoring, classification of answers, or basis for achievement level determination, failing to execute the evaluation task entirely.  \\
  \specialrule{1.2pt}{8pt}{8pt}
  \multicolumn{2}{@{}l}{\textbf{39040}} \\
  \textbf{Question}  & You are designing a multimodal instructional module for an advanced linguistics course that teaches students how to produce persuasive short documentary videos (3--5 minutes) that communicate sociolinguistic research findings to non-specialist audiences. Describe and justify your complete design choices across image, audio, and video modalities. Specifically, (1) identify the target audience and learning objectives; (2) select and justify the multimodal affordances (visual framing, shot types, graphic overlays, color palette, voice style, pacing, ambient sound, music, captions, and image/sound editing techniques) you will use to make complex linguistic concepts accessible without oversimplification; (3) explain how you will structure the 3--5 minute narrative (sequence of scenes, evidence presentation, and persuasive appeals) and tie each segment to specific multimodal decisions; (4) discuss how you will ensure ethical representation of communities and accurate portrayal of data (including consent, anonymization, and visualization choices); and (5) propose a brief evaluation plan to judge whether the video succeeded in communicating the research to the target audience (specify which qualitative and/or quantitative measures you would collect). Justify each decision with reference to multimodal communication principles and principles from linguistics (e.g., indexicality, register, speech accommodation).   \\
  \textbf{Reference} & Target audience and learning objectives: The target audience is educated non-specialists aged 18--35 (undergraduates, policy interns, community organizers) with interest in language and society but limited technical training in linguistics. Learning objectives are: (1) convey the central research question and main empirical findings of a sociolinguistic study in plain language; (2) demonstrate how linguistic evidence supports the claims using clear multimodal examples; (3) foster critical reflection on social implications and ethical considerations. These objectives prioritize comprehension, evidentiary reasoning, and ethical awareness---appropriate for persuasive documentary style aimed at motivating informed action or reflection. \newline Multimodal affordances and justifications: \newline - Visual framing and shot types: Use a mix of medium close-ups for speaker credibility and intimacy (when introducing researchers or community members), wide establishing shots to situate community context, and extreme close-ups for visualizing micro-evidence (e.g., screenshots of transcribed utterances, spectrogram excerpts). Close-ups support ethos and pathos; wide shots provide contextual indexicality showing environment-linked language variation. Macro-to-micro alternation helps viewers move from general claims to concrete evidence. \newline - Graphic overlays and data visualization: Design simple, annotated visuals---color-coded utterance excerpts, timelines showing change over time, and minimally-annotated spectrograms with salient features highlighted. Use consistent visual metaphors (e.g., speech balloons linked to speakers) to map audio examples to text. Avoid dense charts; prefer annotated exemplars that preserve linguistic nuance. This balances accessibility with fidelity to data and leverages dual coding (visual + verbal) to scaffold comprehension. \newline - Color palette and typography: Use a restrained, high-contrast palette (two primary colors plus neutrals) to guide attention without distracting. Typography should use a legible sans-serif for captions and a monospace-like font for transcriptions to signal analytic content. Color and type cues create visual hierarchy and signal transitions between narrative, data, and interpretation. \newline - Voice style and onscreen narration: Adopt a conversational yet precise narratorial voice---moderate speaking rate, warm timbre, and clear enunciation---to maximize comprehensibility and trustworthiness. When presenting technical terms, use brief lay definitions and voice modulation to mark importance. This reflects accommodation: matching register to audience while retaining disciplinary accuracy. \newline - Pacing and editing rhythms: Maintain brisk pacing overall (3--5 minutes) with 8--15 second average shot lengths, lengthening slightly for critical examples to allow cognitive processing. Use rhythmic alternation between claim, example, and interpretation segments so viewers can encode evidence-based reasoning. Cognitive load theory and multimedia learning principles guide pacing choices. \newline - Ambient sound and music: Use unobtrusive ambient sound to situate scenes (market noise, classroom hum) and a low-volume instrumental track to support emotional contour without competing with speech. Silence or sparse texture is used intentionally before key examples to focus attention. Audio layering reinforces contextual indexicality and persuasive tone. \newline - Captions and accessibility: Provide verbatim captions synchronized with speech, indicate paralinguistic features (e.g., laughter, pauses) in brackets for analytic transparency, and offer on-screen transcripts of key utterances with IPA or orthographic forms depending on audience familiarity. Accessibility both broadens reach and preserves analytical detail. \newline \textit{[\ldots\ truncated \ldots]} \newline - Evidence-centered persuasion: Presenting raw exemplars followed by explicit analytical pointers demonstrates linkage from data to claim, modeling critical evaluation rather than mere assertion. \newline - Ethical linguistics: Transparency about transcription choices, anonymization, and community review respects epistemic responsibility and reduces the risk of extractive representation. \newline Together, these design choices ensure the short documentary is persuasive yet rigorous: it makes principled trade-offs between simplicity and fidelity, uses multimodal affordances to scaffold understanding, and embeds ethical safeguards to respect and accurately represent linguistic communities.      \\
  \textbf{Score}     & 2/10  \\
  \textbf{Reasoning} & [1-2] The answer is severely incomplete, cutting off mid-sentence in section 2 and entirely missing sections 3-5. This results in limited learning effectiveness, no evidence of accessibility/ethical considerations, and failure to integrate the required educational design elements.  \\
  \specialrule{1.2pt}{8pt}{8pt}
  \multicolumn{2}{@{}l}{\textbf{41448}} \\
  \textbf{Question}  & You are a teaching assistant monitoring a university-level ensemble rehearsal. During a live performance, a section of student musicians consistently plays a passage one beat late and with incorrect articulation compared to the score. Describe, in detail, how you would detect and analyze the root causes of this error in real time and provide at least three concrete corrective suggestions you would deliver immediately to the instructor to help the students fix the issue during the same rehearsal. In your analysis, break down possible sensory, cognitive, and pedagogical factors (e.g., score reading, ensemble cues, rehearsal technique, tempo perception, notation misinterpretation) and explain how each factor could produce the observed late entries and wrong articulation. For each corrective suggestion, explain why it targets specific root causes, how the instructor should implement it in the moment (exact phrasing or gesture if relevant), and what measurable cue or behavior the instructor should look for to confirm improvement within the rehearsal.   \\
  \textbf{Reference} & Detection and analysis: \newline First, I would confirm the pattern by listening for consistency rather than isolated instances: the passage is one beat late and has an articulation mismatch across multiple measures and multiple players in the same section. I would quickly check the score at that passage to confirm the written articulation and placement of notes relative to the beat (e.g., is there an anacrusis, syncopation, or pickup that could be misread?). I would observe the players' body language and breath/cueing behavior---are they watching the conductor, a section leader, or the prior barline? Are they anticipating a cue or reacting late? I would also note the tempo and conductor gestures: is the tempo steady, or has the conductor's preparatory beat been ambiguous? Finally, I would consider rehearsal history: was that passage recently changed (dynamic/articulation) or has it been rehearsed less than surrounding material? \newline Possible root causes (broken down): \newline - Score-reading/notation misinterpretation: The passage might include syncopation, tied notes, or an upbeat that students misread as a downbeat, causing entries to fall one beat late. Incorrectly interpreted articulation markings (e.g., staccato vs. tenuto) or phrasing slurs could cause uniform articulation errors. \newline - Cueing and ensemble leadership: The section may be waiting for a visual or aural cue from the conductor or another section and misinterpreting the cue timing (e.g., following the end of a previous phrase rather than the intended beat). If a section leader cues late, the rest will follow late. \newline - Tempo perception and subdivision: Students may be hearing larger beats and not subdividing correctly; a fast tempo with weak internal subdivision can push their entries to the next beat. If they're counting in a way that emphasizes the wrong pulse, entries slip by one beat. \newline - Cognitive load and memory: If the passage is technically difficult or recently changed, students may be mentally processing fingerings or rhythms, delaying their entry and causing imprecise articulation. \newline - Articulation misunderstanding or technical limitation: Players may default to a habitual articulation (e.g., legato) conflicting with the written staccato or accents, either from misunderstanding or from focusing on timing rather than articulation mechanics. \newline - Acoustics and balance: If the conductor's beat or a conductor's preparatory breath is not audible to the section (due to room acoustics or distance), they may rely on sound cues that arrive late. \newline \textit{[\ldots\ truncated \ldots]} \newline - Measurable cue of improvement: The section's attacks become visually and aurally synchronized with the designated leader; latency between leader and rest of section drops to near zero on the next take, and articulation becomes uniform with the leader's model. \newline In all three suggestions, provide immediate feedback language the instructor can use: concise, prescriptive prompts such as "Watch my prep---one-and-two-and---play on one,Slow---exaggerate staccato---three times," or "Principal breath---everyone in sync---play." These brief commands reduce ambiguity and cognitive load. \newline Final check and next steps: After implementing one or two of the corrections, the instructor should run the passage twice at tempo and listen for consistent on-beat attacks and uniform articulation. If timing errors persist but articulation improves, focus next on tempo/subdivision drills. If timing clears but articulation remains, do isolated articulation exercises (e.g., repeated-note staccato at metronome). Document whether the cause appears primarily notational misunderstanding, cueing, tempo perception, or technique so the instructor can schedule follow-up micro-rehearsals addressing the root cause.      \\
  \textbf{Score}     & 4/10  \\
  \textbf{Reasoning} & [3-4] The error analysis is generally accurate and well-structured, but the answer is incomplete and cuts off mid-sentence. Crucially, it entirely omits the required three concrete corrective suggestions, implementation steps, and measurable cues, making the response inapplicable for direct use in the lesson context.  \\
  \specialrule{1.2pt}{8pt}{8pt}
  \multicolumn{2}{@{}l}{\textbf{94826}} \\
  \textbf{Question}  & You are an evaluator using an essay rubric to assess a kindergarten student's short narrative (about 4--6 sentences) written in simple English. Analyze the student's essay by: 1) identifying the specific rubric criteria (e.g., organization, vocabulary, sentence structure, conventions, creativity) that are met and those that are not met; 2) citing at least three concrete pieces of textual evidence (short quoted phrases or sentences from the essay) to support each claim; and 3) explaining how each piece of evidence maps to the rubric descriptors (showing the reasoning that connects evidence to criterion). Provide a clear, organized analysis that could be shared with the classroom teacher to justify your scoring decisions.   \\
  \textbf{Reference} & Student essay (example): "I went to the park. I saw a big dog. The dog ran fast. I felt happy and I played with the dog." \newline Analysis: \newline 1) Organization: Partially met. Evidence: "I went to the park." (opening setting), "I saw a big dog." (event follows setting), and "I felt happy and I played with the dog." (concluding feeling/action). These three phrases show a basic sequence: setting event reaction, which aligns with the rubric descriptor for simple narrative organization (clear beginning, middle, end). However, the transitions between ideas are minimal and there is no explicit ending sentence that wraps up the story, so the essay meets the organization criterion at a basic level but not at a more developed level. \newline 2) Vocabulary and Word Choice: Mostly met for age level. Evidence: "big dog,ran fast," and "happy." These concrete, age-appropriate words convey the main images and emotions clearly. The use of descriptive words like "big" and "fast" demonstrates an attempt to add detail. The vocabulary is simple but effective for kindergarten expectations, mapping to a rubric descriptor that values appropriate and concrete word choice for the grade. \newline 3) Sentence Structure and Fluency: Partially met. Evidence: "I went to the park." (complete simple sentence), "I saw a big dog." (complete simple sentence), and "I felt happy and I played with the dog." (compound sentence connected with "and"). The presence of complete sentences shows control of basic sentence formation; the compound sentence indicates emerging ability to combine ideas. However, all sentences are short and formulaic, so while the student demonstrates foundational sentence structure, they do not yet show varied or complex sentences required for higher rubric levels. \newline 4) Conventions (spelling, capitalization, punctuation): Mostly met. Evidence: Each sentence begins with a capital "I" and ends with a period in the sample text provided. Words like "park,dog," and "happy" are spelled correctly. This matches rubric descriptors for correct use of basic conventions at the kindergarten level. One possible area for attention is spacing and any capitalization of proper nouns if present in other drafts, but in this sample conventions are largely correct. \newline 5) Creativity and Elaboration: Minimally met. Evidence: "I saw a big dog" and "I played with the dog" show a simple event, but there is little elaboration about the dog (color, name, actions beyond running) or sensory detail beyond "ran fast" and "felt happy." The briefness of the narrative and lack of imaginative details indicate the student meets only the lowest rubric level for elaboration. \newline Mapping evidence to rubric descriptors (reasoning): \newline - "I went to the park." maps to the organization descriptor for a clear beginning because it establishes setting at the start. \newline \textit{[\ldots\ truncated \ldots]} \newline - Correct periods and capitalization map directly to the conventions descriptor for correct basic mechanics. \newline - The absence of additional sensory details, characterization, or varied sentence types maps to rubric levels indicating minimal elaboration and limited complexity. \newline Overall evaluative synthesis: The essay demonstrates solid foundational skills expected of a kindergarten writer: clear simple sentences, appropriate basic vocabulary, correct basic conventions, and a recognizable narrative sequence (beginning event reaction). Areas for growth include adding more descriptive detail (who the dog was, what it looked like, how it acted) and varying sentence length/structure to increase elaboration and narrative interest. The textual evidence cited above supports awarding middle-level ratings for organization, vocabulary, and conventions, lower-middle ratings for sentence structure, and a low rating for creativity/elaboration.      \\
  \textbf{Score}     & 1/10  \\
  \textbf{Reasoning} & [1-2] The student's answer completely omits the required analysis, evidence, and mapping, making the result meaningless for the evaluation task due to the missing input.  \\
  \specialrule{1.2pt}{8pt}{8pt}
  \multicolumn{2}{@{}l}{\textbf{14308}} \\
  \textbf{Question}  & You are preparing an environmental analysis report for a kindergarten English Language Arts program. Based on your analysis of the user's educational values, compatibility with existing educational systems, and the physical and social characteristics of the facility, identify and prioritize three core educational values that should guide program decisions. For each value: (a) explain why it is prioritized given the user's context and compatibility with existing systems; (b) describe one likely conflict that could arise between this value and another stakeholder value or system constraint; and (c) propose a practical strategy to resolve or mitigate that conflict in the classroom and facility design. Support your reasoning with specific examples tied to kindergarten ELA practices and the facility's features.   \\
  \textbf{Reference} & After reviewing the user's stated educational values and the characteristics of the existing education system and facility, I prioritize these three core educational values for a kindergarten ELA program: child-centered playfulness, language-rich inclusivity, and safe routine consistency. Below I justify each choice, identify likely conflicts, and propose practical resolution strategies grounded in kindergarten ELA practices and the facility's features. 1) Child-centered playfulness --- Priority justification: The user values early childhood engagement and the existing system allows some curricular flexibility; play-based approaches align with developmental ELA outcomes (oral language development, emergent literacy) and are well suited to a kindergarten schedule. The facility's features---if there are open play areas, low shelving, and natural light---support hands-on centers, storytelling corners, and interactive language games. Conflict: A likely conflict is pressure from standardized expectations or rigid curriculum pacing in the broader education system that emphasizes measurable skills over free play. Teachers or administrators may feel required to allocate more time to formal literacy drills, reducing play opportunities. Resolution strategy: Integrate playful activities that intentionally target ELA objectives so play counts as assessment time. For example, convert a dramatic play corner into a "post office" where children write and dictate simple notes (letter formation, vocabulary) and teachers use anecdotal records to document language use. Schedule short, focused explicit instruction blocks (10--15 minutes) for phonological awareness, paired with extended playful centers where those same skills are practiced in authentic contexts. Use the facility's open area for rotating play-based literacy stations to demonstrate alignment with system standards while preserving play-centered pedagogy. 2) Language-rich inclusivity --- Priority justification: The user emphasizes respecting children's home languages and diverse cultural backgrounds; the existing education system may require English outcomes but can accommodate multilingual supports. The facility's social characteristics (multicultural community, family involvement) make inclusivity vital for engagement and equitable language development. Conflict: Tension may arise between inclusive multilingual practices and stakeholders who prioritize rapid English-only proficiency or who lack resources for bilingual materials. Parents or administrators seeking quick measurable English gains may view time spent on home-language use as delaying progress. Resolution strategy: Adopt additive bilingual approaches that validate home languages while scaffolding English. Practically, create dual-language book corners and picture labels in children's home languages plus English, and involve families in creating word walls from home-language story contributions. Use echo-reading and shared book routines where teachers model English while encouraging children to retell stories in their home language, then bridge to English vocabulary. Leverage facility spaces like a family-resource corner for bilingual materials and parent-led story sessions, demonstrating how multilingual practices support, rather than hinder, English acquisition through transfer and increased language confidence. 3) Safe routine consistency --- Priority justification: Young learners thrive on predictable routines for emotional security, which supports affective engagement and readiness for literacy learning. If the facility has limited transitional spaces or a large class size, consistent routines reduce chaos and maximize instructional time. The existing system may require attendance, reporting, and behavior expectations that are better met when children have stable daily structures. Conflict: Strict routines can conflict with the child-centered value of flexibility and exploratory play; staff turnover or scheduling constraints may also make maintaining consistent routines challenging. Resolution strategy: Co-design predictable yet flexible routines with staff and families. For example, establish a consistent morning literacy ritual (welcome song, shared reading, choice time) that always includes a language-rich component but allows children to choose centers that reflect evolving interests. Use visual schedules and classroom signals placed at child eye level to support transitions, and configure facility zones (reading nook, writing table, sensory area) so that even when teacher schedules change, the physical cues maintain routine. To address staff turnover, produce a concise classroom routine guide and short video for substitutes showing how routines embed ELA opportunities (e.g., morning message writing, group storytelling). Concrete examples tying values to practice and facility features: - In a bright corner near windows (facility feature), create a multilingual book nook (value: language-rich inclusivity) where families donate books labeled in home languages and children practice retelling stories during circle time. - Use open floor space and low shelves to set up thematic play centers (value: child-centered playfulness) where children role-play bookstore or puppet theater, integrating vocabulary and print concepts; teachers rotate in to document language use for portfolios aligning with system reporting. - Install pictorial schedules at child height near classroom entry (value: safe routine consistency) so children independently move between shared-reading, choice centers, and dismissal, reducing disruptions and increasing time-on-task for ELA activities. In summary, prioritizing child-centered playfulness, language-rich inclusivity, and safe routine consistency balances the user's educational values with the realities of the existing system and facility. Anticipating conflicts---such as standardized pressures, English-only expectations, and tension between routine and flexibility---and using concrete, classroom-level strategies (integrated playful instruction, additive bilingual resources, visual routines tied to the physical space) will help resolve value conflicts while promoting affective engagement and emergent literacy outcomes.      \\
  \textbf{Score}     & 4/10  \\
  \textbf{Reasoning} & [3-4] The report omits major issues by being incomplete, cutting off mid-sentence in Priority 2 and entirely missing Priority 3. This results in low completeness, fitting the 3-4 band, though the first section shows some logical structure.  \\
  \specialrule{1.2pt}{8pt}{8pt}
  \bottomrule
\end{longtable}